% !LW recipe = latexmk (xelatex)
\documentclass[12pt,a4paper,UTF8]{ctexart}
\usepackage{SJTUReport}
\usepackage{booktabs,tabularx,longtable,mathtools,enumitem,needspace,tikz}
\usepackage[section]{placeins}
\usetikzlibrary{arrows.meta,positioning,shapes.geometric}
\numberwithin{equation}{section}
\setlength{\emergencystretch}{2em}
\setlength{\headheight}{15pt}
\setlist[enumerate]{itemsep=3pt,topsep=5pt,parsep=0pt}
\setlist[itemize]{itemsep=3pt,topsep=5pt,parsep=0pt}
\renewcommand{\arraystretch}{1.18}
\newcommand{\todoresult}[1]{\textcolor{gray}{#1}}
\newcommand{\ind}{\mathbb I}
\newcommand{\R}{\mathcal R}
\def\subsectionautorefname{节}
\title{\vspace{1.6cm}\heiti\LARGE 基于强化学习终端价值的高速公路\\[0.4em]换电路径与充电功率滚动时域优化\par\vspace{1.2cm}}
\author{\relax}
\date{}

\begin{document}
\cover
\thispagestyle{empty}
\clearpage
\begin{abstract}
面向高速公路长途出行场景，本文研究预约与随机两类用户共同参与的换电站运营问题。车辆在一次行程中可能多次换电，换电站的选择决定沿线各站的服务需求和退回电池状态，充电功率则影响电池恢复服务能力的时间及购电费用。本文对预约用户的换电路径和换电站充电功率进行滚动时域优化，利用强化学习估计预测域的终端价值，综合考虑终端库存和预测域外未完成服务的影响。模型根据车辆续航构建候选网络，以日前换电路径作为在线调整的基准，并结合电池 SOC、等待窗口、服务顺序和站点总充电功率上限描述运营过程。充电功率与换电路径采用不同的更新周期，每轮执行当前时间段的决策，再根据新观测更新模型。价值函数通过运营轨迹中的实际收益训练，并以分段线性形式纳入优化目标。数值实验分别围绕换电路径与充电功率联合优化、RL 终端价值和运行条件设置对照，为评价所提方法提供统一的实验方案。

\medskip
\noindent\textbf{关键词：}高速公路；换电站；滚动时域优化；充电功率；强化学习；终端价值
\end{abstract}
\clearpage
\tableofcontents
\clearpage

% !TeX root = ../main.tex
\section{引言}

电动汽车的长途出行需要沿高速公路获得连续、可靠的补能服务。换电站通过交付已经充好的电池缩短车辆的现场补能时间，其服务能力同时取决于车辆到站分布和站内电池库存。长途车辆可能在多个站点换电，前一次选站会改变后续可达站点；每次换电还会向站内退回一块具有特定 SOC 的电池，使车辆的换电安排与各站未来的充电需求相互影响。

运营中同时存在预约用户和随机到站用户。预约信息为需求安排提供依据，车辆实际进入高速公路后的位置、SOC 和预计到站时间又会不断更新。随机用户的到站时间和退回电池状态具有不确定性。与此同时，分时电价、设备功率和站级供电能力共同影响充电安排。因此，运营者需要结合实时信息选择预约用户的换电站点，并协调各站充电功率，以平衡服务收入、购电费用、路径调整和预约履约。

换电路径与充电功率具有直接联系。将车辆安排到某一站点，会改变该站的电池需求、退回 SOC 和后续充电负荷；提高某一站点的充电功率，也可能使原本库存紧张的换电安排获得可行的服务时间。本文将两类决策纳入同一滚动时域优化模型，在每轮预测域内统一评价其运营影响。充电功率按较短周期更新，换电路径按较长周期更新，以适应站内状态变化并保持用户换电安排的稳定。

有限预测域会截断电池库存和预约服务的长期影响。预测结束时仍在充电的电池可能服务于后续需求，已经确定的预约行程也可能在预测域外继续换电。终端价值用于评价这些状态对后续运营的影响。已有研究通过近似价值迭代构造预测控制的终端成本，并分析其与预测长度及价值近似误差的关系\cite{moreno2022terminal}。本文结合高速公路换电场景，将终端库存与预测域外未完成服务共同纳入状态价值函数，通过实际运营轨迹训练，并在每轮优化中计算本轮决策对应的终端价值。

本文的研究内容包括三个方面。首先，利用 SOC 分档候选网络描述预约用户的可达换电路径，建立换电路径与充电功率的联合优化模型。其次，在模型中表达满电交付、退回电池充电、预约优先、等待窗口及路径调整费用，使决策与电池和用户状态的变化相互对应。最后，构造包含终端库存和未完成服务的价值函数，并通过联合优化消融、终端价值消融及敏感性分析评价各组成部分的作用。

图~\ref{fig:framework} 给出整体方法。第 2 节介绍问题与基本假设；第 3 节建立候选网络、日前换电路径和滚动时域优化模型；第 4 节介绍终端价值学习；第 5 节介绍数值实验；第 6 节总结全文。附录给出分段线性函数和相关变量乘积的等价约束。

\begin{figure}[!htbp]
\centering
\begin{tikzpicture}[
  box/.style={draw=blue!45!black,rounded corners=2pt,align=center,
    text width=3.8cm,minimum height=1.1cm,fill=blue!3,font=\small},
  arr/.style={-{Latex[length=2mm]},thick,draw=blue!55!black}]
  \node[box] (obs) {实时状态与预测\\车辆、请求、电池、电价};
  \node[box,right=0.75cm of obs] (opt) {滚动时域优化\\换电路径与充电功率};
  \node[box,right=0.75cm of opt] (act) {当前时间段执行\\服务、充电与状态更新};
  \node[box,below=1.1cm of opt] (value) {RL 终端价值\\终端库存与未完成服务};
  \draw[arr] (obs)--(opt);
  \draw[arr] (opt)--(act);
  \draw[arr] (value)--node[right,font=\small]{价值函数}(opt);
  \draw[arr] (act.south)|-node[pos=0.3,right,font=\small]{实际收益与轨迹}(value.east);
  \draw[arr] (act.north)--++(0,0.6)-|node[pos=0.5,above,font=\small]{下一时刻观测}(obs.north);
\end{tikzpicture}
\caption{换电路径与充电功率滚动时域优化方法}
\label{fig:framework}
\end{figure}


% !TeX root = ../main.tex
\section{问题描述与基本假设}

\subsection{系统与运营场景}
考虑由多个 O--D 对及沿线换电站组成的高速公路网络。车辆从入口进入后沿下游方向行驶，经必要的换电站补能并从出口离开。各站配备固定数量的电池和充电设备。一次换电取走满电电池并退回带有剩余电量的电池，退回电池经过充电后继续服务其他用户。因此，换电站的选择既影响车辆的后续行程，也影响各站电池 SOC 和充电需求。

系统服务预约用户和随机用户。预约用户提供 O--D、计划入口时间和入口 SOC，研究中的预约用户均具有满足续航条件的完整换电路径。尚未进入高速公路的用户由预约信息描述；进入高速公路后，使用其位置、SOC 和预计到站时间更新换电安排。已经到站且正在等待的请求保留原到站时刻、截止时刻和退回 SOC。随机用户直接到站提出请求，未来随机需求由当前预测信息描述。

\subsection{基本假设}
为描述上述运营过程，作如下假设。
\begin{enumerate}[label=\textbf{假设\arabic*.},leftmargin=4.4em]
  \item \textbf{道路与行驶。}车辆沿高速公路下游方向行驶，并按照给定换电站序完成行程。节点间距离和行驶能耗为已知参数，换电路径满足续航和出口最低 SOC 要求。研究场景中的入口 SOC 与初始路径相容，给定行驶能耗下已确定的换电安排在两次路径更新之间保持续航可行。
  \item \textbf{电池规格与库存。}可交换电池具有统一容量和兼容规格，服务就绪 SOC 归一化为 $1$。每次换电取走和退回各一块电池，站内电池数量保持不变。SOC $1$ 表示运营规定的服务目标。
  \item \textbf{充电过程。}充电效率为已知参数，电池 SOC 按输入电量线性变化。不同站点可以具有不同充电效率，各充电设备和站点分别具有给定功率上限。
  \item \textbf{换电操作。}换电视为瞬时操作。交付电池达到服务就绪 SOC，退回电池进入对应槽位并参加后续充电。槽位记录当时接入的电池状态。
  \item \textbf{信息与预测。}每轮可获得当前电池状态、等待请求及在途车辆的位置和 SOC。车辆到站时间采用当前 ETA 预测，未来随机需求和电价采用决策时刻可得的信息。同一用户沿下游方向的 ETA 按行驶顺序递增，预测信息随滚动过程更新。
  \item \textbf{到站顺序。}同一站点同类请求的原始到站时刻互不相同。同一时间段内可以存在多个请求，其先后次序由原始到站时刻确定。
\end{enumerate}

\subsection{运营规则}
预约请求优先于随机请求，同类请求按原到站时刻先后服务。每个请求自到站时刻起具有给定的最大等待时间，运营者可以在该窗口内选择服务时间，并为其匹配可用满电电池。预约在截止时刻仍未获得服务时，产生一次预约失败成本，该用户后续换电需求随之结束；随机请求在等待截止后流失。已完成的换电和行驶过程按实际记录保留，在线路径调整针对车辆后续的换电安排。换电安排实际发生变化时产生路径调整成本。


% !TeX root = ../main.tex
\section{模型建立}

除特别说明外，用户集合、请求集合及可用路径网络均指当前优化轮对应的集合，其轮次索引予以省略。各集合随滚动过程更新，用户与请求的标识保持不变。

将运营时间划分为长度为 $\delta$ 的时间段，记
\begin{subequations}\label{eq:time_grid}
\begin{align}
t_n&=n\delta,\label{eq:time_points}\\
\mathcal I_n&=[t_n,t_{n+1}),\label{eq:time_intervals}\\
\mathcal T_\ell&=\{\ell,\ldots,\ell+H-1\}.\label{eq:prediction_indices}
\end{align}
\end{subequations}
其中 $\ell$ 为当前轮，$H$ 为预测时间段数。本轮预测范围为 $[t_\ell,t_{\ell+H})$，在 $t_n$ 处理服务及换电，随后在 $\mathcal I_n$ 内充电。每轮只执行当前时间段，并在下一时刻使用新观测重新优化。默认 $\delta=5$~min，能量计算时取 $\delta=1/12$~h。

充电功率每段更新，换电路径每 $K$ 段允许更新一次。以运营起点为共同基准，定义已知参数
\begin{subequations}\label{eq:update_clock}
\begin{align}
g_\ell&=\ind\{\ell\bmod K=0\},\label{eq:path_update_indicator}\\
K&=3.\label{eq:path_update_interval}
\end{align}
\end{subequations}
因此，两类预约用户的路径均每 $K\delta=15$~min 允许更新一次。$g_\ell=1$ 时联合优化全部预约用户的路径与充电功率；$g_\ell=0$ 时，在途用户保持最近发布的换电安排，尚未进入高速公路的用户保持上一次优化后保留的计划，仅优化充电功率。未进入高速的用户首次参与优化时，以日前路径作为初始计划。进入高速前，优化后的计划不向用户发布，也不产生实际调整费用。图~\ref{fig:updates} 示意两种决策的更新安排。

\begin{figure}[!htbp]
\centering
\begin{tikzpicture}[x=2.55cm,y=0.7cm,font=\small]
\draw[-{Latex[length=2mm]},thick] (0,0)--(5.15,0);
\foreach \x/\lab in {0/0,1/5,2/10,3/15,4/20,5/25}{
 \draw (\x,-0.08)--(\x,0.08);
 \node[above] at (\x,0.13) {\lab~min};
 \fill[teal!70!black] (\x,-0.7) circle (2.3pt);
}
\foreach \x in {0,3}{\fill[orange!90!black] (\x,-1.5) rectangle ++(0.10,0.16);}
\node[anchor=east] at (-0.14,-0.7) {充电};
\node[anchor=east] at (-0.14,-1.4) {路径};
\node[below,align=center] at (2.5,-1.9) {每轮更新车辆、电池和请求状态，并执行当前时间段决策};
\end{tikzpicture}
\caption{充电功率与换电路径的更新安排}
\label{fig:updates}
\end{figure}

表~\ref{tab:notation} 按集合与索引、输入参数、决策变量、由决策计算的量与价值函数列出主要符号。预约用户按 O--D 对分组，以 $\mathcal P$ 表示 O--D 对集合，$p\in\mathcal P$ 为组索引，$k$ 为该组内的用户索引。用户由 $(p,k)$ 唯一标识，编号在滚动过程中保持不变。到站请求统一以 $r$ 标识；某一请求所在站点记为 $i(r)$。费用和价格参数统一采用小写 $c$，合计费用采用大写 $C$，服务收入记为 $I$。

\begin{longtable}{@{}p{4.4cm}p{11.1cm}@{}}
\caption{主要符号}\label{tab:notation}\\
\toprule 符号 & 含义 \\ \midrule
\endfirsthead
\caption[]{主要符号（续）}\\
\toprule 符号 & 含义 \\ \midrule
\endhead
\bottomrule\endfoot
\multicolumn{2}{@{}l@{}}{\textbf{集合与索引}}\\*
$\mathcal S;\ i,j$ & 换电站集合及站点索引 \\
$\mathcal B_i;\ b$ & 站点 $i$ 的电池充电槽集合及充电槽索引 \\
$\mathcal P;\ p$ & O--D 对集合及索引 \\
$\mathcal K^{\mathrm{fut},p},\mathcal K^{\mathrm{enr},p}$ & O--D 对 $p$ 下本轮尚未进入、已经进入高速公路的预约用户集合 \\
$\mathcal K^p;\ k$ & O--D 对 $p$ 下本轮预约用户集合及组内用户索引, $\mathcal K^p=\mathcal K^{\mathrm{fut},p}\cup\mathcal K^{\mathrm{enr},p}$ \\
$\R,\R_i;\ r$ & 本轮全部请求集合、站点 $i$ 的请求集合及请求索引 \\
$\R_i^A,\R_i^R$ & 站点 $i$ 的预约请求集合与随机请求集合 \\
$\mathcal V^{p,k},\mathcal A^{p,k}$ & 用户 $(p,k)$ 本轮路径网络的节点集合与可用弧集 \\
$\mathcal T_\ell;\ n,\ell$ & 本轮预测时间段集合；$n$ 为时段索引，$\ell$ 为当前轮索引 \\
\midrule
\multicolumn{2}{@{}l@{}}{\textbf{输入参数}}\\*
$\delta$ & 时间段长度 \\
$H,K$ & 预测域长度与路径更新间隔，均以时间段数计 \\
$o^{p,k},d^p$ & 用户 $(p,k)$ 本轮路径的起点与出口节点 \\
$v_{i,j}^p$ & O--D 对 $p$ 中节点 $i$ 至 $j$ 的行驶 SOC 消耗 \\
$E_B$ & 电池容量（kWh） \\
$\eta_i$ & 站点 $i$ 的充电效率 \\
$\rho_r$ & 请求 $r$ 对应的退回电池 SOC \\
$S_{i,b}^{\mathrm{obs}}$ & 本轮开始时充电槽 $b$ 中电池的观测 SOC \\
$\overline P_{i,b}^{\mathrm{ch}},\overline P_i^{\mathrm{sta}}$ & 设备充电功率上限与站点总充电功率上限（kW） \\
$\bar t_r^{\mathrm{arr}}$ & 当前等待请求的实际到站时刻，或不需要先在本轮完成另一次换电的请求的预计到站时刻 \\
$\tau_{j,i}^{p,k}$ & 用户 $(p,k)$ 从换电站 $j$ 到下一换电站 $i$ 的预测行驶时间，本轮求解时已知且为正 \\
$\Delta$ & 请求自到站起允许的最大等待时间 \\
$\bar y_{i,j}^{p,k}$ & 用户 $(p,k)$ 的日前换电路径是否使用弧 $(i,j)$ \\
$c_{i,n}^{\mathrm{ch}},c_{i,n}^{\mathrm{sw}}$ & 充电电价与换电服务单价（元/kWh） \\
$c^{\mathrm{adj}},c^{\mathrm{fail}}$ & 单次路径调整惩罚与单次预约失败成本（元/次） \\
$\bar x_i^{p,k},\widehat x_i^{p,k},x_i^{p,k,\mathrm{pub}}$ & 本轮开始时已知的日前路径、保留计划和最近发布路径对应的换电站安排 \\
$\bar{\boldsymbol c}_r,\boldsymbol c_{r,m}$ & 本轮求解前根据车辆行驶预测计算的候选终端车辆信息 \\
$\beta$ & 终端价值权重\\
\midrule
\multicolumn{2}{@{}l@{}}{\textbf{决策变量}}\\*
$y_{i,j}^{p,k}$ & 用户 $(p,k)$ 的本轮换电路径是否经过弧 $(i,j)$，二元变量 \\
$P_{i,b,n}$ & 站点 $i$ 的设备 $b$ 在时间段 $n$ 的充电功率（kW） \\
$\alpha_{b,r,n}$ & 是否在 $t_n$ 用充电槽 $b$ 的电池服务请求 $r$，二元变量 \\
\midrule
\multicolumn{2}{@{}l@{}}{\textbf{由决策计算的量与价值函数}}\\*
$S_{i,b,n}^{-}$ & 充电槽 $b$ 中电池在 $t_n$ 当次换电前的 SOC，由初始电量、换电安排和充电功率确定 \\
$\chi_{p,k}$ & 本轮的换电站安排是否与日前路径（未进入高速）或最近发布路径（在途）不同，由路径比较确定 \\
$f_r$ & 预约请求 $r$ 是否在本轮预测期内导致用户失败，由是否需要服务、截止时刻和换电次数确定 \\
$Q_{r,n}$ & 请求 $r$ 在 $t_n$ 是否已到站且仍在等待换电，由到站时间和此前换电次数确定 \\
$w_r$ & 请求 $r$ 在预测终点是否仍需服务，由本轮换电和失败情况确定 \\
$t_r^{\mathrm{arr}},t_r^{\mathrm{ddl}}$ & 请求的到站时刻与等待截止时刻；后站到站时刻随前站换电安排确定 \\
$B_r$ & 在本轮预测终点是否能够确定请求 $r$ 的到站时刻，由前站换电安排确定 \\
$e_{r,n},h_r$ & 分别表示 $t_n$ 是否处于到站至截止时间内，以及截止时刻是否早于预测终点；到站时间尚未确定时均为 $0$ \\
$Z_r(n),Z_r$ & 分别表示请求 $r$ 在本轮预测中 $t_n$ 之前的累计换电次数，以及本轮预测期内的累计换电次数。 \\
$I$ & 本轮预测期内的服务收入（元） \\
$C^{\mathrm{ch}},C^{\mathrm{adj}},C^{\mathrm{fail}}$ & 本轮预测期内的充电费用、本轮路径调整惩罚及本轮预测期内的预约失败费用（元） \\
$s_\ell,s_{\ell+H}$ & 本轮观测运营状态与由本轮决策确定的预测终端状态 \\
$V_\theta(s)$ & 强化学习估计的状态终端价值（元） \\
\end{longtable}

\subsection{基于 SOC 分档的候选网络}
\label{sec:network}
对每个 O--D 对 $p$，以入口、沿线换电站和出口构造网络。车辆路径
$o^p\to i_1\to\cdots\to i_q\to d^p$
表示依次在 $i_1,\ldots,i_q$ 换电，弧所跨越的其他站点不被访问。将入口 SOC 划分为若干区间，分别预先生成候选网络，以减少在线需要处理的路径组合。

节点集合为 $\mathcal V^p=\{o^p\}\cup\mathcal S^p\cup\{d^p\}$，所有节点沿下游方向排列。以 $D_{i,j}^p$ 和 $v_{i,j}^p$ 表示节点间距离和行驶 SOC 消耗，以 $\underline s_d^p$ 表示出口 SOC 下限。最小换电间距 $\underline D$ 用于优先排除距离过近的首段和站间换电安排，出口末段按续航条件处理。对每个 SOC 区间 $h$，按以下四步生成候选网络。
\begin{enumerate}[label=\textbf{步骤\arabic*.},leftmargin=4.2em]
\item \textbf{生成原始弧。}若以该 SOC 区间上限电量能够从入口到达下游站点 $j$，则连接入口与 $j$；若 $v_{i,j}^p\le1$，则连接下游站点 $i,j$；若 $v_{i,d^p}^p+\underline s_d^p\le1$，则连接站点与出口。以该 SOC 区间上限电量从入口直接驶向出口后仍满足出口 SOC 下限时，加入相应弧。
\item \textbf{生成完整路径。}在原始网络中枚举所有入口至出口的完整路径。
\item \textbf{剪枝过近换电弧。}找出距离小于 $\underline D$ 的首段弧和站间弧，按距离从小到大考察。若当前仍存在至少一条不包含该弧的完整路径，则删除该弧；否则保留。
\item \textbf{形成候选网络。}合并最终保留路径中的弧，得到 SOC 区间对应的弧集 $\mathcal A^{h,p}$。
\end{enumerate}

在线使用时，根据用户实际入口 SOC 进一步检查首段续航。对在途用户，删除已经经过的节点，以实时位置建立虚拟起点，按实时 SOC 检查其到下游站点的首段弧。最小换电间距始终以入口或最近一次实际换电位置为基准。对已经在站等待的用户，将当前站作为后续路径起点，后续行程以该次换电完成后的满电状态计算；当前等待请求在第~\ref{sec:model}~节中保留，并与后续服务关联。

从本轮起点 $o^{p,k}$ 直接驶向出口 $d^p$ 的弧，只有满足以下相应条件时才予以保留：
\begin{equation}
\begin{cases}
s_{\mathrm{dep}}^{p,k}-v_{o^{p,k},d^p}^{p,k}\ge\underline s_d^p,
&\text{用户正在行驶或尚未进入高速},\\
1-v_{o^{p,k},d^p}^{p}\ge\underline s_d^p,
&\text{用户当前在站等待换电}.
\end{cases}
\label{eq:direct_exit_soc}
\end{equation}
式~\eqref{eq:direct_exit_soc} 第一行使用本轮起点的已知车辆 SOC $s_{\mathrm{dep}}^{p,k}$，其中 $v_{o^{p,k},d^p}^{p,k}$ 为从该起点直接驶向出口的 SOC 消耗。第二行使用当前站换电后的满电状态。两种情况都要求车辆到达出口时的 SOC 不低于 $\underline s_d^p$。由此得到用户弧集 $\mathcal A^{p,k}$、起点 $o^{p,k}$ 以及对应出口 $d^p$。

\subsection{日前换电路径计算}
\label{sec:dayahead}
日前换电路径根据预约用户的 O--D、入口 SOC 和候选网络计算。先使用实际预约入口 SOC 检查首段弧，并保留满足出口 SOC 要求的完整路径。设所得完整路径集合为 $\mathcal Y^{p,k}$，其元素由弧指示量 $y$ 表示，则日前路径取为
\begin{equation}
\bar y^{p,k}\in\arg\min_{y\in\mathcal Y^{p,k}}
\sum_{(j,i)\in\mathcal A^{p,k}:\,i\in\mathcal S}y_{j,i}.
\label{eq:dayahead_path}
\end{equation}
该目标选择换电次数最少的完整路径。当多条路径的换电次数相同时，依次比较其第一、第二及后续换电站的位置，优先选择更靠近出口的站点；站点位置相同的记录使用固定网络顺序区分。得到的 $\bar y^{p,k}$ 及对应换电站序列作为在线路径调整的基准。该计算仅需要车辆与道路信息，输入用户均具有满足上述条件的完整可达路径。

\subsection{滚动时域优化模型}
\label{sec:model}

\paragraph{请求与输入信息。}
对于用户 $k\in\mathcal K^p$，对每条到达换电站 $i$ 的可用弧 $(j,i)\in\mathcal A^{p,k}$，建立请求 $r=(p,k,j,i)$，其站点为 $i(r)=i$，退回 SOC 为
\begin{equation}
\rho_r=
\begin{cases}
s_{\mathrm{dep}}^{p,k}-v_{o^{p,k},i}^{p,k},&j=o^{p,k},\ \text{用户正在行驶或尚未进入高速},\\
1-v_{j,i}^{p},&j\text{ 为换电站}.
\end{cases}
\label{eq:return_soc}
\end{equation}
其中 $v_{o^{p,k},i}^{p,k}$ 为本轮起点到站点 $i$ 的已知 SOC 消耗，$s_{\mathrm{dep}}^{p,k}$ 为本轮起点的已知车辆 SOC，未到用户使用预约信息及入口前观测，在途用户使用实时值。用户当前在站等待时，后续首段按第二种情况计算。各记录的 $\rho_r$ 在求解前已知，并由第~\ref{sec:network}~节的网络生成规则保证 $0\le\rho_r<1$。

令 $\R_i^A,\R_i^R$ 分别为站点 $i$ 的预约请求集合和随机请求集合，$\R_i=\R_i^A\cup\R_i^R$，并以 $\R^A,\R^R,\R$ 表示各站相应集合的并集。

以 $\R_{p,k}^A$ 表示预约用户 $(p,k)$ 在本轮模型中的所有换电请求。这些请求包括当前在站等待的换电请求，以及后续行程中可能产生的换电请求。本轮可用弧对应的预约请求组成集合
\begin{equation}
\begin{aligned}
\R^{\mathrm{path}}=\bigl\{(p,k,j,i)\in\R:\;&
p\in\mathcal P,\ k\in\mathcal K^p,\\
&(j,i)\in\mathcal A^{p,k},\ i\in\mathcal S\bigr\}.
\end{aligned}
\label{eq:path_request_set}
\end{equation}
式~\eqref{eq:path_request_set} 包含所有预约用户沿本轮可用弧可能产生的换电请求。对于已经在站等待的预约用户，当前站是本轮路径的起点，驶入该站的弧不在本轮可用弧集合中。该用户当前等待的请求单独保留，后续沿可用弧产生的请求仍属于 $\R^{\mathrm{path}}$。因此，$\R\setminus\R^{\mathrm{path}}$ 由当前在站等待的预约请求和全部随机请求组成。

为描述同一用户的换电先后顺序，定义集合
\begin{equation}
\mathcal P_r=
\begin{cases}
\bigl\{q\in\R_{p,k}^A:i(q)=j\bigr\},
&r=(p,k,j,i)\in\R^{\mathrm{path}},\\
\varnothing,&r\in\R\setminus\R^{\mathrm{path}},
\end{cases}
\qquad r\in\R.
\label{eq:predecessor_request_set}
\end{equation}
其中，$q$ 表示用户 $(p,k)$ 的一个换电请求，$i(q)$ 为该请求所在的站点。式~\eqref{eq:predecessor_request_set} 第一行将该用户在前一站 $j$ 的所有可能换电请求放入 $\mathcal P_r$。用户可能通过不同路径到达 $j$，但所选路径至多采用其中一个请求。若用户当前已在 $j$ 等待换电，则从 $j$ 前往下一换电站的请求，其 $\mathcal P_r$ 只包含这次等待请求。

当前等待请求自身和随机请求按式~\eqref{eq:predecessor_request_set} 第二行取空集。对于从入口或当前行驶位置出发的首次换电请求，本轮起点没有尚需服务的换电请求，第一行得到的集合也为空。$\mathcal P_r=\varnothing$ 表示在本轮安排请求 $r$ 换电之前，不需要先完成该用户的另一次换电。本轮开始前已经完成的换电不再包含在请求集合中。上述集合均在求解前确定。

以 $a_r$ 表示是否需要为请求 $r$ 安排换电。对于后续可能发生的换电请求 $r=(p,k,j,i)$，所选路径经过弧 $(j,i)$ 时取 $a_r=1$，否则取 $a_r=0$，即 $a_r=y_{j,i}^{p,k}$。当前在站等待的请求和随机请求不受预约用户路径选择的影响，均取 $a_r=1$。$a_r=1$ 表示有这项换电需求，是否完成换电还取决于后续服务约束。

当前在站等待的请求使用实际到站时刻。对于用户当前未在站等待时的下一次换电，以及随机请求，预计到站时刻由本轮观测和预测给定。上述时刻记为 $\bar t_r^{\mathrm{arr}}$，在本轮求解中保持不变。已经到站的请求保留原始到站时刻；尚未到站的首次请求和随机请求在下一轮根据新观测更新预测。

同一用户后续站点的到站时刻由前一站的换电时刻与站间行驶时间共同确定。以 $\tau_{j,i}^{p,k}>0$ 表示用户 $(p,k)$ 从换电站 $j$ 到下一换电站 $i$ 的预测行驶时间。该参数在本轮求解前给定，求解过程中保持不变，下一轮根据新观测更新。前一站等待时间的变化通过换电时刻传递到后一站。到站时间、等待截止时间及其与服务安排的关系在下文约束中定义。

\paragraph{决策变量与计算量。}
本轮优化选择预约用户的换电路径、各请求的换电时刻和所用电池，以及各充电槽的充电功率。用二元变量 $y_{i,j}^{p,k}\in\{0,1\}$ 表示用户 $(p,k)$ 的换电路径是否经过弧 $(i,j)$。一组满足流平衡的 $y$ 变量唯一确定该用户依次访问的换电站。用非负连续变量 $P_{i,b,n}\ge0$ 表示站点 $i$ 的充电槽 $b$ 在时间段 $n$ 内的充电功率。

以 $\alpha_{b,r,n}\in\{0,1\}$ 表示是否在 $t_n$ 用站点 $i(r)$ 的充电槽 $b$ 中的电池服务请求 $r$，其中 $b\in\mathcal B_{i(r)}$。电池在 $t_n$ 当次换电前的 SOC 记为 $S_{i,b,n}^{-}$，由初始电量、换电安排和充电功率递推得到。

下文的路径改变标记 $\chi_{p,k}$、预约失败标记 $f_r$、等待标记 $Q_{r,n}$ 和终端未完成服务标记 $w_r$ 均由相应定义确定，取值为 $0$ 或 $1$。$\chi_{p,k}$ 比较本轮选择的换电站与已有路径：未进入高速的用户与日前路径比较，在途用户与最近发布的路径比较。

$f_r=1$ 表示请求 $r$ 需要服务、等待截止时刻早于本轮预测终点，且仍未完成换电，从而导致该用户预约失败。$Q_{r,n}=1$ 表示请求在 $t_n$ 已到站、尚未超过等待截止时刻，并且此前尚未换电。$w_r=1$ 表示预测终点仍需继续处理该请求。同一用户失败后，其后续请求不再保留。

本轮开始时的日前路径、保留计划和最近发布路径均为已知输入。其余计算量在对应约束中定义；求解时所需的辅助变量和线性化按附录~\ref{lin:appendix} 处理。

\paragraph{目标函数。}

\begin{equation}
\max\quad I-C^{\mathrm{ch}}-C^{\mathrm{adj}}-C^{\mathrm{fail}}
+\beta V_\theta(s_{\ell+H}),
\label{eq:objective}
\end{equation}
其中，各项分别为
\begin{subequations}\label{eq:objective_terms}
\begin{align}
I&=E_B\sum_{n\in\mathcal T_\ell}\sum_{r\in\R}
c_{i(r),n}^{\mathrm{sw}}(1-\rho_r)\sum_{b\in\mathcal B_{i(r)}}\alpha_{b,r,n},\label{eq:income}\\
C^{\mathrm{ch}}&=\delta\sum_{n\in\mathcal T_\ell}\sum_{i\in\mathcal S}
c_{i,n}^{\mathrm{ch}}\sum_{b\in\mathcal B_i}P_{i,b,n},\label{eq:charging_cost}\\
C^{\mathrm{adj}}&=c^{\mathrm{adj}}\sum_{p\in\mathcal P}\sum_{k\in\mathcal K^p}\chi_{p,k},\label{eq:adjustment_cost}\\
C^{\mathrm{fail}}&=c^{\mathrm{fail}}\sum_{r\in\R^A}f_r.\label{eq:failure_cost}
\end{align}
\end{subequations}

式~\eqref{eq:income} 中，$I$ 为本轮预测期内的换电服务收入。每次换电交付满电电池，并收回 SOC 为 $\rho_r$ 的电池，因此按补充电量 $E_B(1-\rho_r)$ 与换电时刻对应的服务单价计费。

式~\eqref{eq:charging_cost} 中，$C^{\mathrm{ch}}$ 为本轮预测期内的充电费用。其中，$c_{i,n}^{\mathrm{ch}}$ 表示站点 $i$ 在时间段 $n$ 的充电电价，充电槽 $b$ 在时间段 $n$ 内从电网输入的电量为 $\delta P_{i,b,n}$。

式~\eqref{eq:adjustment_cost} 中，$C^{\mathrm{adj}}$ 为本轮优化中的路径调整惩罚。尚未进入高速公路的用户与日前路径比较，在途用户与最近一次向其发布的路径比较。只要新增或取消任意一个换电站，就对该用户计一次 $c^{\mathrm{adj}}$，改变多个站点也只计一次。

未进入高速的用户偏离日前路径时，该惩罚只用于本轮优化，不计入实际运营费用，也不在多轮之间累计结算。实际调整费用只在向在途用户发布不同于原安排的路径时计入，具体见式~\eqref{eq:actual_reward}。

式~\eqref{eq:failure_cost} 中，$C^{\mathrm{fail}}$ 为本轮预测期内的预约失败费用。当预约请求因超过等待截止时刻仍未完成换电而导致预约失败时计入一次费用 $c^{\mathrm{fail}}$。

$V_\theta(s_{\ell+H})$ 估计从预测终点状态出发的后续累计净收益。其输入 $s_{\ell+H}$ 由式~\eqref{eq:terminal_state} 定义，价值函数在第~\ref{sec:value}~节的式~\eqref{eq:value_network} 中给出。它考虑终点时的电池电量和未完成换电请求，使本轮决策也考虑预测范围之外的充电与服务需求。$\beta$ 调整后续净收益在目标函数中的权重。


\paragraph{（1）流平衡约束。}
对每个用户，路径满足
\begin{equation}
\sum_{j:(i,j)\in\mathcal A^{p,k}}y_{i,j}^{p,k}
-\sum_{j:(j,i)\in\mathcal A^{p,k}}y_{j,i}^{p,k}
=\begin{cases}
1,&i=o^{p,k},\\-1,&i=d^p,\\0,&\text{其他节点},
\end{cases}
\quad i\in\mathcal V^{p,k},\ p\in\mathcal P,\ k\in\mathcal K^p.
\label{eq:flow}
\end{equation}
式~\eqref{eq:flow} 确定一条起点至出口的路径。

\paragraph{（2）路径调整约束。}
令 $x_i^{p,k}=\sum_{j:(j,i)\in\mathcal A^{p,k}}y_{j,i}^{p,k}$ 表示预约用户 $(p,k)$ 是否在站点 $i$ 换电。$\bar x_i^{p,k}$ 表示日前路径是否安排该用户在该站换电。对于尚未进入高速公路的用户，以 $\widehat x_i^{p,k}$ 表示本轮开始时保留的计划是否安排其在该站换电；首次参与优化时取 $\widehat x_i^{p,k}=\bar x_i^{p,k}$，以后使用上一轮保留的计划。对于在途用户，$x_i^{p,k,\mathrm{pub}}$ 表示最近一次向其发布的路径是否安排其在该站换电。已经完成的换电和当前正在等待的换电不参与路径调整，比较各条路径时同时去除这些站点；当前等待请求单独保留。各路径未安排换电的站点取 $0$。

两类用户都只能在式~\eqref{eq:update_clock} 规定的时刻调整路径，因此对各 $i\in\mathcal S$、$p\in\mathcal P$，有
\begin{equation}
\begin{aligned}
|x_i^{p,k}-\widehat x_i^{p,k}|&\le g_\ell,
&&k\in\mathcal K^{\mathrm{fut},p},\\
|x_i^{p,k}-x_i^{p,k,\mathrm{pub}}|&\le g_\ell,
&&k\in\mathcal K^{\mathrm{enr},p}.
\end{aligned}
\label{eq:path_freeze}
\end{equation}
式~\eqref{eq:path_freeze} 在 $g_\ell=0$ 时要求未进入高速的用户保持已保留的计划，在途用户保持最近发布的路径。$g_\ell=1$ 时，两类用户都可以重新选择换电站。

路径调整惩罚按以下关系计算：
\begin{equation}
\chi_{p,k}=\begin{cases}
\displaystyle\max_{i\in\mathcal S}|x_i^{p,k}-\bar x_i^{p,k}|,
&k\in\mathcal K^{\mathrm{fut},p},\\
\displaystyle\max_{i\in\mathcal S}|x_i^{p,k}-x_i^{p,k,\mathrm{pub}}|,
&k\in\mathcal K^{\mathrm{enr},p},
\end{cases}
\qquad p\in\mathcal P.
\label{eq:path_change}
\end{equation}
式~\eqref{eq:path_change} 中，任意一个站点的换电安排与比较路径不同，就有 $\chi_{p,k}=1$。未进入高速的用户可以根据供需变化提前调整计划，但在进入高速前不发布，以避免反复通知用户。惩罚始终相对日前计划计算，以固定的比较依据限制不必要的路径调整。式~\eqref{eq:path_freeze} 和式~\eqref{eq:path_change} 的线性化方法见附录~\ref{lin:appendix}。

\paragraph{（3）换电次数约束。}
式~\eqref{eq:path_request_set} 中的 $\R^{\mathrm{path}}$ 表示本轮可用弧对应的预约请求，其余请求为当前在站等待的预约请求和随机请求。用 $a_r$ 表示是否需要为请求 $r$ 安排换电，其取值为
\begin{equation}
a_r=\begin{cases}
y_{j,i}^{p,k},&r=(p,k,j,i)\in\R^{\mathrm{path}},\\
1,&r\in\R\setminus\R^{\mathrm{path}},
\end{cases}
\qquad r\in\R.
\label{eq:request_selection}
\end{equation}
式~\eqref{eq:request_selection} 第一行规定：所选路径经过弧 $(j,i)$ 时，$a_r=1$；否则 $a_r=0$。第二行将当前在站等待的请求和随机请求的 $a_r$ 固定为 $1$。因此，$a_r$ 的取值由上述等式确定；取 $1$ 时能否完成换电，还需满足后续服务约束。

请求 $r$ 在 $t_n$ 之前的累计换电次数定义为
\begin{equation}
Z_r(n)=\sum_{\substack{m\in\mathcal T_\ell\\m<n}}
\sum_{b\in\mathcal B_{i(r)}}\alpha_{b,r,m}.
\label{eq:service_sums}
\end{equation}
式~\eqref{eq:service_sums} 适用于 $n\in\mathcal T_\ell\cup\{\ell+H\}$，只计算从 $t_\ell$ 到 $t_n$ 之前完成的换电，不包含 $t_n$ 当次的换电。为简化记号，记 $Z_r=Z_r(\ell+H)$，表示请求 $r$ 在本轮预测期内的累计换电次数。

每个请求在本轮预测期内最多完成一次换电，同一充电槽在同一时刻最多交付一块电池。$\alpha_{b,r,n}=1$ 表示在 $t_n$ 用充电槽 $b$ 中的电池为请求 $r$ 换电。相应约束为
\begin{subequations}\label{eq:assignment_capacity}
\begin{align}
Z_r&\le a_r\quad(r\in\R),\label{eq:request_service_limit}\\
\sum_{r\in\R_i}\alpha_{b,r,n}&\le1
\quad(i\in\mathcal S,b\in\mathcal B_i,n\in\mathcal T_\ell).
\label{eq:battery_service_limit}
\end{align}
\end{subequations}
式~\eqref{eq:request_service_limit} 规定：$a_r=0$ 时，该请求不能获得服务；$a_r=1$ 时，最多获得一次服务。式~\eqref{eq:battery_service_limit} 将同一充电槽在 $t_n$ 服务的请求数限制为最多一个。其中，$\R$ 为本轮全部请求集合，$\R_i$ 为站点 $i$ 的请求集合，$\mathcal B_i$ 为该站的充电槽集合；约束对所有站点 $i\in\mathcal S$ 和本轮各时刻 $t_n$（$n\in\mathcal T_\ell$）成立。

\paragraph{（4）到站时刻与等待期限。}
用户到站后才能换电，超过等待截止时刻后不再安排服务。对于需要在多个站点换电的用户，下一站的到站时刻取决于前一站的换电时刻。若前一站在本轮预测期内没有换电，下一站的到站时间暂时不能确定。

用 $B_r$ 表示在本轮预测终点是否能够确定请求 $r$ 的到站时刻，定义
\begin{equation}
B_r=\begin{cases}
1,&\mathcal P_r=\varnothing,\\
\displaystyle\sum_{q\in\mathcal P_r}Z_q,&\mathcal P_r\ne\varnothing,
\end{cases}
\qquad r\in\R.
\label{eq:predecessor}
\end{equation}
其中，$\mathcal P_r$ 是同一用户在前一站的所有可能换电请求，见式~\eqref{eq:predecessor_request_set}。$\mathcal P_r=\varnothing$ 时，到站时刻由输入给定，所以 $B_r=1$。否则，式~\eqref{eq:predecessor} 对前一站请求在本轮预测期内的换电次数求和。所选路径至多采用其中一个请求，因此前一站完成换电时 $B_r=1$，未完成时 $B_r=0$。

$t_r^{\mathrm{arr}}$ 和 $t_r^{\mathrm{ddl}}$ 分别表示请求 $r$ 的到站时刻和等待截止时刻。若用户 $(p,k)$ 需要先在站点 $j$ 换电，再前往站点 $i$，则到达 $i$ 的时刻等于在 $j$ 的换电时刻加上该用户的预测行驶时间 $\tau_{j,i}^{p,k}$。对所有 $r\in\R$，有
\begin{equation}
\begin{aligned}
t_r^{\mathrm{arr}}&=\begin{cases}
\bar t_r^{\mathrm{arr}},&\mathcal P_r=\varnothing,\\
\displaystyle\sum_{q\in\mathcal P_r}\sum_{m\in\mathcal T_\ell}
(t_m+\tau_{j,i}^{p,k})\sum_{b\in\mathcal B_j}\alpha_{b,q,m},
&\mathcal P_r\ne\varnothing,
\end{cases}\\
t_r^{\mathrm{ddl}}&=t_r^{\mathrm{arr}}+\Delta B_r.
\end{aligned}
\label{eq:arrival_time}
\end{equation}
式~\eqref{eq:arrival_time} 第一种情况使用给定的实际或预计到站时刻。第二种情况使用本轮安排的前站换电时刻。例如，前一站在 10:25 换电、站间行驶需 30~min，则下一站在 10:55 到站，等待截止时刻再加上 $\Delta$。

$B_r=0$ 时，式~\eqref{eq:arrival_time} 将到站和截止时刻都填为 $0$，表示尚不能确定这两个时刻。该请求只要属于所选路径且用户尚未失败，仍作为未完成需求保留，由终端价值评价后续服务。

用 $e_{r,n}$ 表示 $t_n$ 是否位于请求到站之后、等待截止之前，包含两个端点。对所有 $r\in\R$ 和未来时刻 $n\in\{\ell+1,\ldots,\ell+H\}$，定义
\begin{equation}
e_{r,n}=B_r\ind\{t_r^{\mathrm{arr}}\le t_n\le t_r^{\mathrm{ddl}}\}.
\label{eq:service_window}
\end{equation}
式~\eqref{eq:service_window} 在到站时间已确定且时间条件满足时取 $1$，其余情况取 $0$。$B_r=0$ 时不会把填零的时刻误当作实际到站或截止时刻。$e_{r,n}=1$ 只说明时间条件满足，能否换电还要看其他约束。到站时刻由本轮服务安排确定，因此未来的 $e_{r,n}$ 也随这些决策变化。

当前的 $e_{r,\ell}$ 由真实到站和等待情况确定：只有已经在站等待且尚未超过等待截止时刻的请求取 $1$。因此，当前只服务实际已经到站的请求；未来时刻按式~\eqref{eq:service_window} 计算，预测终点只计算状态。

\paragraph{（5）预约失败与随机流失。}
用 $h_r$ 表示请求的等待截止时刻是否早于本轮预测终点，定义
\begin{equation}
h_r=B_r\ind\{t_r^{\mathrm{ddl}}<t_{\ell+H}\},
\qquad r\in\R.
\label{eq:deadline_passed}
\end{equation}
式~\eqref{eq:deadline_passed} 只比较截止时刻与预测终点，不判断请求是否已经获得服务，也不比较服务时刻与截止时刻。例如，预测终点为 11:00、截止时刻为 10:30，即使请求已经在 10:20 换电，仍有 $h_r=1$。到站时间尚不能确定时，$B_r=0$，相应的 $h_r=0$。截止恰好等于预测终点时也取 $0$，因为该请求仍有在该时刻换电的机会。当前实际在站等待的请求均未超过截止时刻；所有使用给定到站时刻的请求满足 $\bar t_r^{\mathrm{arr}}+\Delta\ge t_\ell$。

预约失败标记由
\begin{equation}
f_r=a_rh_r(1-Z_r),\qquad r\in\R^A
\label{eq:failure_exact}
\end{equation}
确定。式~\eqref{eq:failure_exact} 表示：只有请求需要服务、截止时刻早于预测终点，且仍未完成换电时，才取 $f_r=1$。已经按时完成换电的请求，即使 $h_r=1$，也因 $Z_r=1$ 而有 $f_r=0$。式~\eqref{eq:queue_status} 和式~\eqref{eq:queue_service} 将换电限制在 $e_{r,n}=1$ 的时刻，因此不允许在等待截止之后服务。

每位预约用户在本轮预测期内至多失败一次，记
\begin{equation}
F_{p,k}=\sum_{r\in\R_{p,k}^A}f_r\le1,
\qquad p\in\mathcal P,\ k\in\mathcal K^p.
\label{eq:failure_order}
\end{equation}
$F_{p,k}$ 表示该用户是否在本轮预测期内失败，用于在终端状态中取消其后续需求。前一站失败意味着该站没有完成换电，后续站点的到站时刻因而尚未确定，相应的 $h_r$ 和 $f_r$ 均为 $0$，失败费用只计算一次。

截止不早于预测终点时，不能因为本轮尚未安排换电就提前判为失败。仍需服务的请求保留在终端状态，由后续滚动优化继续安排。随机请求的流失量为 $h_r(1-Z_r)$，即截止早于预测终点且仍未服务时计为流失。

\paragraph{（6）等待和服务顺序。}
本组确定每个时刻哪些请求正在等待换电，并规定多人同时等待时的服务优先顺序。等待标记和服务安排满足
\begin{subequations}\label{eq:queue}
\begin{equation}
Q_{r,n}=a_re_{r,n}[1-Z_r(n)],
\qquad r\in\R,\quad n\in\mathcal T_\ell\cup\{\ell+H\}.
\label{eq:queue_status}
\end{equation}
\begin{equation}
\sum_{b\in\mathcal B_{i(r)}}\alpha_{b,r,n}\le Q_{r,n},
\qquad r\in\R,\quad n\in\mathcal T_\ell.
\label{eq:queue_service}
\end{equation}
\end{subequations}
式~\eqref{eq:queue_status} 表示：请求需要服务、已经到站且尚未超过截止时刻，并且此前尚未换电时，才取 $Q_{r,n}=1$。$Z_r(n)$ 不包含 $t_n$ 当次换电，所以本次获得服务的请求在服务前仍有 $Q_{r,n}=1$。式~\eqref{eq:queue_service} 只允许为正在等待的请求换电，但不要求每个等待请求都立即获得服务。预测终点只计算 $Q_{r,\ell+H}$，不安排当次换电。

对于需要先在前一站换电的请求，式~\eqref{eq:arrival_time} 与式~\eqref{eq:service_window} 已保证：只有前一站换电并经过正的行驶时间后，$e_{r,n}$ 才能取 $1$。因此，前站尚未换电时不能服务后站，同一用户也不能在同一时刻完成前后两次换电。若某站预约在 $t_n$ 之前已经失败，该请求因超过截止时刻而有 $e_{r,n}=0$。在该用户所选路径上，后续请求因前站未换电而有 $e_{r,n}=0$，更早的请求已经服务，$Z_r(n)=1$；未选择的请求则有 $a_r=0$。这些条件已使失败用户的请求退出等待，无需另设逐时刻的用户失败标记。上述乘积关系均为准确的定义，取值为 $0$ 或 $1$，线性化按附录~\ref{lin:appendix} 处理。

在同一站点定义 $q\prec r$：当 $q$ 为预约而 $r$ 为随机请求，或者二者同类且 $t_q^{\mathrm{arr}}<t_r^{\mathrm{arr}}$ 时成立。同类请求的先后顺序由式~\eqref{eq:arrival_time} 中本轮确定的到站时刻比较得到。对任意两个不同请求，施加以下条件约束：
\begin{equation}
\begin{gathered}
q\prec r\quad\Longrightarrow\quad
\sum_{b\in\mathcal B_i}\alpha_{b,r,n}
\le 1-Q_{q,n}+\sum_{b\in\mathcal B_i}\alpha_{b,q,n},\\
q,r\in\R_i,\quad q\ne r,\quad i\in\mathcal S,\quad n\in\mathcal T_\ell.
\end{gathered}
\label{eq:priority}
\end{equation}
式~\eqref{eq:priority} 要求：较高优先级请求仍在等待时，只有它在同一时刻获得服务，才能服务较低优先级请求。尚未到站、已经服务或已经失效的请求，其 $Q_{q,n}=0$，不会阻止其他请求服务。同类请求同时到站时不规定二者的先后。每个到站时刻只可能取给定时刻、有限个 $t_m+\tau_{j,i}^{p,k}$ 或未确定时的零值，因此上述严格先后关系可按有限个已知时刻的组合精确展开。运营者可以在等待期限内选择换电时刻，并通过 $\alpha_{b,r,n}$ 确定各请求使用的电池。

\paragraph{（7）满电交付和充电。}
在每个时刻 $t_n$，先换电，再在随后的时间段 $\mathcal I_n$ 内充电。以 $S_{i,b,n}^{-}$ 表示充电槽 $b$ 中电池在 $t_n$ 当次换电前的 SOC，将换电与随后充电的影响直接写入下一时刻的电量：
\begin{subequations}\label{eq:battery_dynamics}
\begin{align}
S_{i,b,\ell}^{-}&=S_{i,b}^{\mathrm{obs}},\label{eq:initial_soc}\\
S_{i,b,n}^{-}&\ge\alpha_{b,r,n},
&&r\in\R_i,\label{eq:full_delivery}\\
S_{i,b,n+1}^{-}&=S_{i,b,n}^{-}
-\sum_{r\in\R_i}(1-\rho_r)\alpha_{b,r,n}
+\frac{\eta_i\delta}{E_B}P_{i,b,n},\label{eq:charge_soc}\\
0\le S_{i,b,n}^{-}&\le1,
&&n\in\mathcal T_\ell\cup\{\ell+H\},\label{eq:soc_bounds}\\
0\le P_{i,b,n}&\le\overline P_{i,b}^{\mathrm{ch}},\label{eq:device_power}\\
\sum_{b\in\mathcal B_i}P_{i,b,n}&\le\overline P_i^{\mathrm{sta}}.\label{eq:station_power}
\end{align}
\end{subequations}
上述关系对所有 $i\in\mathcal S$ 和 $b\in\mathcal B_i$ 成立；站点总功率约束对各站点成立。除初始时刻和式~\eqref{eq:soc_bounds} 另有规定外，时间索引均取 $n\in\mathcal T_\ell$。

式~\eqref{eq:initial_soc} 将初始 SOC 固定为当前观测值。式~\eqref{eq:full_delivery} 保证满电交付：$\alpha_{b,r,n}=1$ 时，结合式~\eqref{eq:soc_bounds} 得到 $S_{i,b,n}^{-}=1$；$\alpha_{b,r,n}=0$ 时，不要求该电池满电，也不要求满电电池必须交付。

式~\eqref{eq:charge_soc} 中，下一时刻的 SOC 等于当前 SOC，减去换电造成的下降量，再加上本时段充电带来的增量。由式~\eqref{eq:battery_service_limit}，同一充电槽在同一时刻最多服务一个请求。若服务请求 $r$，满电电池被取走，退回电池的 SOC 为 $\rho_r$，因此扣减 $1-\rho_r$；若没有换电，求和项为 $0$。由于 $0\le\rho_r<1$，换电后槽内电池的 SOC 自然位于 $0$ 与 $1$ 之间。

充电功率 $P_{i,b,n}$ 是从电网输入的功率，$\delta P_{i,b,n}$ 为输入电量，乘以充电效率 $\eta_i$ 后得到存入电池的电量，再除以容量 $E_B$，得到 SOC 增量。功率以 kW 计、容量以 kWh 计，因此 $\delta$ 以小时计。例如，容量为 $60$~kWh、效率为 $0.9$、功率为 $40$~kW 时，充电 $5$~min 使 SOC 增加 $0.05$，即 $5$ 个百分点；若换电后电池的 SOC 为 $0.30$，下一时刻就是 $0.35$。

式~\eqref{eq:soc_bounds} 覆盖本轮各换电时刻和预测终点，保证电量非负且不超过容量。终点的 SOC 已包含最后一个时间段的充电结果，因此同样需要这一上界。

式~\eqref{eq:device_power} 限制单个充电槽的功率，取 $0$ 表示该时间段不充电。所选功率在该时间段内保持恒定，并与 SOC 上界共同限制充电量。式~\eqref{eq:station_power} 限制同一站点各充电槽的总功率。例如，两台设备的功率上限各为 $60$~kW，而站点上限为 $80$~kW 时，两台设备不能同时以 $60$~kW 充电。每个请求可以选择任意满足条件的充电槽，模型同时确定其后续充电功率。

\paragraph{（8）终端状态。}
目标函数~\eqref{eq:objective} 中的 $V_\theta(s_{\ell+H})$ 根据本轮预测终点的状态评价后续净收益。本组确定终点时的电池电量、未完成请求及其等待情况、用户的车辆与路径信息，并将它们汇总为 $s_{\ell+H}$。这些信息如何编码以及价值函数如何计算，在第~\ref{sec:value}~节说明。

终端电池状态取 $S_{i,b,\ell+H}^{-}$，由式~\eqref{eq:battery_dynamics} 递推得到。它包含本轮最后一个时间段的充电结果，不包含预测终点当时的换电，用于评价后续满电电池的供应能力和充电需求。

用 $w_r=1$ 表示请求 $r$ 在预测终点仍需继续安排换电。对预约请求，定义
\begin{equation}
w_r=a_r(1-Z_r)(1-F_{p,k}),\qquad r\in\R_{p,k}^A.
\label{eq:terminal_pending}
\end{equation}
式~\eqref{eq:terminal_pending} 表示：请求需要服务、尚未完成换电，并且该用户没有失败时，才保留该请求。因此，前一站尚未换电而到站时间尚不能确定的后续预约，以及预计在本轮预测期之后到站的预约，只要仍需服务，都保留在终端状态中。同一用户一旦失败，其后续请求均不再保留。

对随机请求，取
\begin{equation}
w_r=(1-h_r)(1-Z_r),\qquad r\in\R^R.
\label{eq:terminal_random}
\end{equation}
式~\eqref{eq:terminal_random} 表示：尚未完成换电且等待截止时刻不早于预测终点的随机请求，仍需继续服务。已经完成换电或已经流失的随机请求取 $w_r=0$。

在式~\eqref{eq:queue_status} 中取 $n=\ell+H$，得到终端等待标记 $Q_{r,\ell+H}$。$w_r=1$ 表示后续还需服务，$Q_{r,\ell+H}=1$ 则进一步表示该请求在预测终点已经到站、尚未超过等待截止时刻，并且仍在等待换电。两者描述同一个请求，不重复计为两项需求。对于保留的请求，还使用 $B_r$、$t_r^{\mathrm{arr}}$ 和 $t_r^{\mathrm{ddl}}$ 说明到站与等待期限：$B_r=1$ 时，两个时刻由式~\eqref{eq:arrival_time} 确定；取 $0$ 时，两个时刻填零，表示到站时间尚不能确定。

对每位预约用户，保留 $w_r=1$ 的请求所在站点，并按道路下游方向排列，即得到预测终点尚需完成的换电安排，其中包括仍在等待的当前站点。

以 $\boldsymbol c_{p,k}(\tau)$ 表示时刻 $\tau$ 的车辆位置、车辆 SOC、是否已进入高速公路，以及计算最小换电间距所用的入口或最近一次实际换电位置。这些信息用于判断后续仍可选择哪些换电站，真实状态使用当前观测。对于预测终端，先根据给定的车辆行驶预测规则计算以下候选值。若 $\mathcal P_r=\varnothing$，以 $\bar{\boldsymbol c}_r$ 表示车辆从本轮观测状态驶向请求 $r$ 所在站点、到站后保持等待时的终端车辆信息；若当前已在该站等待，则保持当前车辆位置和 SOC。计算换电间距的基准位置保持本轮观测值。若 $\mathcal P_r\ne\varnothing$，以 $\boldsymbol c_{r,m}$ 表示车辆在 $t_m$ 于前一站完成换电后，以满电状态驶向请求 $r$ 所在站点、到站后保持等待时的终端车辆信息。此时计算换电间距的基准位置更新为前一站 $i(q)=j$。这些候选值是本轮已知输入，由对应用户的预测行驶时间、沿途能耗和初始车辆信息在求解前计算。预测终端取
\begin{equation}
\begin{aligned}
\boldsymbol c_{p,k}(t_{\ell+H})={}&
\sum_{\substack{r\in\R_{p,k}^A\\\mathcal P_r=\varnothing}}
w_r\bar{\boldsymbol c}_r\\
&+\sum_{\substack{r\in\R_{p,k}^A\\\mathcal P_r\ne\varnothing}}
w_r\sum_{m\in\mathcal T_\ell}\sum_{q\in\mathcal P_r}
\sum_{b\in\mathcal B_{i(q)}}\alpha_{b,q,m}\boldsymbol c_{r,m}.
\end{aligned}
\label{eq:terminal_vehicle_state}
\end{equation}
式~\eqref{eq:terminal_vehicle_state} 只保留所选路径中最前面的未完成请求对应的候选值。该请求之前的换电已经完成，其后的请求则因前一站尚未完成换电而没有对应的服务项。因此，只要用户仍有未完成服务，上式恰有一个候选值被选中；本轮换电时刻改变时，终端车辆位置和 SOC 也随之改变。全部服务完成或用户失败时，所有 $w_r=0$，车辆信息填零，该用户不再作为仍需服务的用户保留。

对于仍需服务的预约用户，终端状态同时保留上述车辆信息、日前换电路径、最近发布的剩余路径及当前保留的计划。当前计划采用本轮选择的路径，已完成的站点从剩余安排中移除。日前路径保留原值，发布记录按前述路径发布规则保留；尚未发布的计划与已发布路径分别保存。这些信息用于评价后续可行的换电安排及可能产生的路径调整费用。

以 $\xi(\tau)$ 表示时刻 $\tau$ 的时间与预测信息，包括当前时刻、当前及后续的充电电价和换电服务单价、随机需求预测，以及路径更新相位。预测终点的路径更新相位为 $(\ell+H)\bmod K$。这些信息用于判断后续需求、充电与服务的收益和费用，以及下一次允许调整路径的时刻。

将上述信息汇总，目标函数~\eqref{eq:objective} 使用的终端状态定义为
\begin{equation}
\begin{aligned}
s_{\ell+H}=\Bigl(&
\bigl\{S_{i,b,\ell+H}^{-}\bigr\}_{i\in\mathcal S,\,b\in\mathcal B_i},\\
&
\bigl\{(r,Q_{r,\ell+H},B_r,
t_r^{\mathrm{arr}},t_r^{\mathrm{ddl}},\rho_r)\bigr\}_{r\in\R:\,w_r=1},\\
&
\bigl\{(\boldsymbol c_{p,k}(t_{\ell+H}),\text{日前路径},\text{已发布路径},\text{保留计划})\bigr\}_{
\substack{p\in\mathcal P,\,k\in\mathcal K^p\\
\sum_{r\in\R_{p,k}^A}w_r>0}},\quad
\xi(t_{\ell+H})\Bigr).
\end{aligned}
\label{eq:terminal_state}
\end{equation}
式~\eqref{eq:terminal_state} 的四部分依次为各充电槽的电池 SOC、未完成请求、仍需服务的预约用户信息，以及时间与预测信息。请求 $r$ 连同其预约或随机类型、所属用户、所在站点及前一换电站等属性一并保留；站点位置、各用户的预测站间行驶时间和退回 SOC 沿用本轮给定输入。$w_r$ 决定保留哪些请求，同一用户的多个请求通过 $(p,k)$ 关联。终端状态中的电量、保留请求、到站时刻、车辆状态和换电安排随本轮决策确定，因此 $V_\theta(s_{\ell+H})$ 能够评价这些决策对后续服务的影响。

式~\eqref{eq:flow}--\eqref{eq:terminal_state}、变量范围和第~\ref{sec:value}~节的价值函数共同构成优化问题。到站、等待和截止条件由给定时刻及离散服务安排精确确定；服务先后关系按有限个候选到站时刻展开。再将绝对值、最大值、价值计算及有界变量乘积按附录~\ref{lin:appendix} 转换为线性约束，得到混合整数线性优化模型。

\paragraph{滚动执行与记账。}
在 $t_\ell$ 读取真实在站请求、电池状态、车辆位置、保留的计划和最新预测，建立本轮模型。路径更新时刻同时优化两类预约用户的路径。未进入高速的用户只保存优化后的计划，用于后续预测和下一轮的 $\widehat x_i^{p,k}$，不发布也不计入实际费用。在途用户的路径若与最近发布的安排不同，则发布新路径、计入一次实际调整费用，并将新路径作为后续比较的依据。非路径更新时刻，两类用户都保持各自已有的安排。

按模型确定的请求和电池安排执行当前换电后，在 $\mathcal I_\ell$ 内使用当前充电功率。区间内到站请求在下一时刻参与换电安排，原到站时刻用于计算等待时间。截止位于 $\mathcal I_\ell$ 内且尚未完成换电的请求在该时段判定超时；恰在下一时刻截止的请求留到下一轮先参与换电安排。

下一轮电池状态由实际服务和充电更新，真实等待请求保留原标识及截止时间。车辆位置、下一次换电的预计到站时刻和站间行驶时间根据新观测更新；仍需先完成前站换电的后续请求继续由式~\eqref{eq:arrival_time} 计算到站时刻。预测请求在实际到站后使用真实到站时刻，状态中始终区分已观测请求与预测请求。若上游预约实际失败，该用户后续服务一起结束。新的当前状态成为下一轮初始条件，新的路径更新参数按式~\eqref{eq:update_clock} 计算。

% !TeX root = ../main.tex
\section{终端价值学习}
\label{sec:value}

强化学习用于估计预测域的终端价值，重点考虑终端库存和预测域外未完成服务。相同的电池总电量可能对应不同的满电供应能力；相同的库存也可能面对不同的预约到站时间、等待窗口和后续换电需求。价值函数将这些信息共同映射为一个终端评价，并通过实际运营轨迹中的累计净收益进行训练。

\subsection{状态与终端价值}
\label{sec:value_state}

以 $s(\tau)$ 表示时刻 $\tau$ 的完整运营状态，包括电池 SOC、真实等待请求、未完成预约及其换电安排，以及时间、电价和需求预测。预约用户的信息还包括是否已进入高速公路。对未进入高速的用户，保存日前路径和最近优化后保留的未发布计划；对在途用户，保存最近发布的剩余路径以及位置、SOC 和下一换电站预计到站时刻等观测。本节先说明本轮决策如何确定预测终端状态，再给出状态编码和价值函数。

\paragraph{预测终端状态。}
对于从 $t_\ell$ 开始的一轮优化，沿用第~\ref{sec:model}~节中的本轮用户、请求和时间集合。目标函数~\eqref{eq:objective} 使用的终端状态 $s_{\ell+H}$ 由本轮路径、服务和充电决策确定，包括终点时的电池电量、未完成请求及其等待情况，以及用户的车辆与路径信息。

终端电池状态取 $S_{i,b,\ell+H}^{-}$，由式~\eqref{eq:battery_dynamics} 递推得到。它包含本轮最后一个时间段的充电结果，不包含预测终点当时的换电，用于评价后续满电电池的供应能力和充电需求。

用 $w_r=1$ 表示请求 $r$ 在预测终点仍需继续安排换电。对预约用户 $(p,k)$ 的请求，定义
\begin{equation}
w_r=a_r\left(
1-\sum_{n\in\mathcal T_\ell}
\sum_{b\in\mathcal B_{i(r)}}\alpha_{b,r,n}
\right)\left(1-\sum_{q\in\R_{p,k}^A}f_q\right),\qquad r\in\R_{p,k}^A.
\label{eq:terminal_pending}
\end{equation}
其中，$f_q$ 由式~\eqref{eq:failure_exact} 确定，$\sum_{q\in\R_{p,k}^A}f_q$ 为该用户在本轮预测期内超过等待期限仍未获服务的预约请求数。前述约束已保证每位用户至多有一个这样的请求，因此该求和只取 $0$ 或 $1$。式~\eqref{eq:terminal_pending} 表示：请求需要服务、尚未完成换电，并且该用户尚无超过等待期限仍未获服务的预约请求时，才保留该请求。因此，前一站尚未换电而到站时间尚不能确定的后续预约，以及预计在本轮预测期之后到站的预约，只要仍需服务，都保留在终端状态中。同一用户一旦有预约请求超过等待期限仍未获服务，其后续请求均不再保留。

对随机请求，取
\begin{equation}
w_r=(1-h_r)\left(
1-\sum_{n\in\mathcal T_\ell}
\sum_{b\in\mathcal B_{i(r)}}\alpha_{b,r,n}
\right),\qquad r\in\R^R.
\label{eq:terminal_random}
\end{equation}
式~\eqref{eq:terminal_random} 表示：尚未完成换电且等待截止时刻不早于预测终点的随机请求，仍需继续服务。已经完成换电或已经流失的随机请求取 $w_r=0$。

在式~\eqref{eq:queue_status} 中取 $n=\ell+H$，得到终端等待标记 $Q_{r,\ell+H}$。$w_r=1$ 表示后续还需服务，$Q_{r,\ell+H}=1$ 则进一步表示该请求在预测终点已经到站、尚未超过等待截止时刻，并且仍在等待换电。两者描述同一个请求，不重复计为两项需求。随机请求和每位预约用户下一次换电请求的到站时刻由式~\eqref{eq:arrival_eta} 给出，即使预计到站时间在本轮预测期之后，也保留这一时间信息。等待截止时刻可由已确定的到站时刻加 $\Delta$ 得到，无需单独保留。更后面的预约请求继续保留站点、行驶时间和退回 SOC 等信息，其具体到站时间暂不确定。

对每位预约用户，保留 $w_r=1$ 的请求所在站点，并按道路下游方向排列，即得到预测终点尚需完成的换电安排，其中包括仍在等待的当前站点。最靠前的请求对应该用户下一次换电，其后各请求无需在本轮确定具体换电时刻。

以 $\boldsymbol c_{p,k}(\tau)$ 表示时刻 $\tau$ 的车辆位置、车辆 SOC、是否已进入高速公路，以及计算最小换电间距所用的入口或最近一次实际换电位置。这些信息用于判断后续仍可选择哪些换电站，真实状态使用当前观测。对于预测终端，先根据给定的车辆行驶预测规则计算以下候选值。若 $\mathcal P_r=\varnothing$，以 $\bar{\boldsymbol c}_r$ 表示车辆从本轮观测状态驶向请求 $r$ 所在站点、到站后保持等待时的终端车辆信息；若当前已在该站等待，则保持当前车辆位置和 SOC。计算换电间距的基准位置保持本轮观测值。若 $\mathcal P_r\ne\varnothing$，以 $\boldsymbol c_{r,m}$ 表示车辆在 $t_m$ 于前一站完成换电后，以满电状态驶向请求 $r$ 所在站点、到站后保持等待时的终端车辆信息。此时计算换电间距的基准位置更新为前一站 $i(q)=j$。这些候选值是本轮已知输入，由对应用户的预测行驶时间、沿途能耗和初始车辆信息在求解前计算。预测终端取
\begin{equation}
\begin{aligned}
\boldsymbol c_{p,k}(t_{\ell+H})={}&
\sum_{\substack{r\in\R_{p,k}^A\\\mathcal P_r=\varnothing}}
w_r\bar{\boldsymbol c}_r\\
&+\sum_{\substack{r\in\R_{p,k}^A\\\mathcal P_r\ne\varnothing}}
w_r\sum_{m\in\mathcal T_\ell}\sum_{q\in\mathcal P_r}
\sum_{b\in\mathcal B_{i(q)}}\alpha_{b,q,m}\boldsymbol c_{r,m}.
\end{aligned}
\label{eq:terminal_vehicle_state}
\end{equation}
式~\eqref{eq:terminal_vehicle_state} 只保留所选路径中最前面的未完成请求对应的候选值。该请求之前的换电已经完成，其后的请求则因前一站尚未完成换电而没有对应的服务项。因此，只要用户仍有未完成服务，上式恰有一个候选值被选中；本轮换电时刻改变时，终端车辆位置和 SOC 也随之改变。全部服务完成，或该用户已有预约请求超过等待期限仍未获服务时，所有 $w_r=0$，车辆信息填零，该用户不再作为仍需服务的用户保留。

对于仍需服务的预约用户，终端状态同时保留上述车辆信息、日前换电路径、最近发布的剩余路径及当前保留的计划。当前计划采用本轮选择的路径，已完成的站点从剩余安排中移除。日前路径保留原值，发布记录按第~\ref{sec:model}~节的路径发布规则保留；尚未发布的计划与已发布路径分别保存。这些信息用于评价后续可行的换电安排及可能产生的路径调整费用。

以 $\xi(\tau)$ 表示时刻 $\tau$ 的时间与预测信息，包括当前时刻、当前及后续的充电电价和换电服务单价、随机需求预测，以及路径更新相位。预测终点的路径更新相位为 $(\ell+H)\bmod K$。这些信息用于判断后续需求、充电与服务的收益和费用，以及下一次允许调整路径的时刻。

将上述信息汇总，目标函数~\eqref{eq:objective} 使用的终端状态定义为
\begin{equation}
\begin{aligned}
s_{\ell+H}=\Bigl(&
\bigl\{S_{i,b,\ell+H}^{-}\bigr\}_{i\in\mathcal S,\,b\in\mathcal B_i},\\
&
\bigl\{(r,Q_{r,\ell+H},
t_r^{\mathrm{arr}},\rho_r)\bigr\}_{r\in\R:\,w_r=1},\\
&
\bigl\{(\boldsymbol c_{p,k}(t_{\ell+H}),\text{日前路径},\text{已发布路径},\text{保留计划})\bigr\}_{
\substack{p\in\mathcal P,\,k\in\mathcal K^p\\
\sum_{r\in\R_{p,k}^A}w_r>0}},\quad
\xi(t_{\ell+H})\Bigr).
\end{aligned}
\label{eq:terminal_state}
\end{equation}
式~\eqref{eq:terminal_state} 的四部分依次为各充电槽的电池 SOC、未完成请求、仍需服务的预约用户信息，以及时间与预测信息。请求 $r$ 连同其预约或随机类型、所属用户、所在站点及前一换电站等属性一并保留；站点位置、各用户的预测站间行驶时间和退回 SOC 沿用本轮给定输入。$w_r$ 决定保留哪些请求，同一用户的多个请求通过 $(p,k)$ 关联。终端状态中的电量、保留请求、到站时刻、车辆状态和换电安排随本轮决策确定，因此 $V_\theta(s_{\ell+H})$ 能够评价这些决策对后续服务的影响。终端请求和车辆状态中变量乘积的线性化见附录~\ref{lin:terminal_state_section}。

\paragraph{状态编码与价值计算。}
对任意状态时刻 $\tau$ 进行编码时，用户和请求集合均取该时刻的集合，$\mathcal K^p$ 表示该时刻 O--D 对 $p$ 下的预约用户集合，省略其时间标记。

为适应不同数量的请求，采用共享分段线性函数和求和生成固定维度表示。对任意状态时刻 $\tau$，令 $\nu_r(\tau)$ 汇集请求站点和前一换电站的位置、该用户在两站之间的预测行驶时间、用户 O--D、退回 SOC、下述时间特征，以及该请求是否正在等待。不存在前一换电站的请求，相应的站点和行驶时间分量填零。预约请求的信息还包含对应用户的车辆信息 $\boldsymbol c_{p,k}(\tau)$、日前路径、最近发布路径和当前保留的计划。预测终端尚需完成的换电站由 $w_r=1$ 的请求确定，车辆状态按式~\eqref{eq:terminal_vehicle_state} 计算，同一用户的请求使用相同的车辆信息。保留的计划与最近发布的路径分别记录，尚未发布的优化结果不替代已发布路径。

以 $w_r(\tau)$ 表示该状态下请求仍需服务，以 $Q_r(\tau)$ 表示其中已经到站且正在等待的请求。真实状态取观测和现有预约，预测终端分别取 $w_r$ 和 $Q_{r,\ell+H}$。对每位预约用户，将仍需服务的请求按道路下游方向排列，最靠前的一个就是下一次换电请求，其中包括当前在站等待的请求。

随机请求和用户下一次换电请求的时间信息使用到站、截止时刻相对状态时刻 $\tau$ 的时间差。真实状态中，已到站请求使用实际到站时刻，尚未到站的下一次请求和随机需求使用给定预计到站时刻。其余后续预约暂不确定具体时间，两个时间分量均填零，继续保留站点、预测行驶时间、退回 SOC 和所属用户等信息。

在预测终点，请求特征 $\nu_r(t_{\ell+H})$ 中的两个时间分量直接按前站换电安排计算。下式中的 $\mathcal P_r$、$\mathcal T_\ell$ 和服务变量仍沿用本轮优化模型的定义，已经完成换电的前站请求也参与求和：
\begin{equation}
\begin{cases}
\begin{pmatrix}
\bar t_r^{\mathrm{arr}}-t_{\ell+H}\\
\bar t_r^{\mathrm{arr}}+\Delta-t_{\ell+H}
\end{pmatrix},
&\mathcal P_r=\varnothing,\\[2mm]
\displaystyle
\sum_{q\in\mathcal P_r}\sum_{m\in\mathcal T_\ell}
\left(\sum_{b\in\mathcal B_j}\alpha_{b,q,m}\right)
\begin{pmatrix}
t_m+\tau_{j,i}^{p,k}-t_{\ell+H}\\
t_m+\tau_{j,i}^{p,k}+\Delta-t_{\ell+H}
\end{pmatrix},
&\mathcal P_r\ne\varnothing.
\end{cases}
\label{eq:request_time_features}
\end{equation}
式~\eqref{eq:request_time_features} 第一种情况直接使用给定时间。第二种情况中，各候选时间差在本轮求解前已知，前站服务变量选出本轮预测采用的一对时间差；前站没有换电时，求和自然为零。对于仍需服务的预约请求，只有该用户最靠前的请求不再需要等待另一次换电，其时间分量采用实际或预计值；更后面的请求取零。因而真实状态和预测终端使用相同的时间信息。零填充值不表示用户此时到站或已经超时；正在等待的请求，其到站时间差可以为负，截止时间差为零时仍有当时换电的机会。

时间与预测信息 $\xi(\tau)$ 沿用式~\eqref{eq:terminal_state} 前的定义。状态编码为
\begin{equation}
\begin{aligned}
f_B(\tau)&=\mathop{\mathrm{concat}}_{i\in\mathcal S}
\left(\sum_{b\in\mathcal B_i}\phi_B(S_{i,b}^{-}(\tau))\right),\\
f_A(\tau)&=\sum_{p\in\mathcal P}\sum_{k\in\mathcal K^p}
\left[\phi_A\!\left(\sum_{r\in\R_{p,k}^A}w_r(\tau)\psi_A\bigl(\nu_r(\tau)\bigr)\right)-\phi_A(0)\right],\\
f_R(\tau)&=\sum_{r\in\R^R}w_r(\tau)\psi_R\bigl(\nu_r(\tau)\bigr),\\
\Phi_\theta(s(\tau))&=\bigl[f_B(\tau),f_A(\tau),f_R(\tau),\xi(\tau)\bigr].
\end{aligned}
\label{eq:state_encoding}
\end{equation}
其中 $\phi_B,\phi_A,\psi_A,\psi_R$ 为固定维度的分段线性映射，使用仿射变换和 ReLU 构造。站内电池使用共享映射，以反映 SOC 分布；预约先按用户汇总再整体编码，以保留同一行程多次换电的关联。对每位用户减去 $\phi_A(0)$，使已经没有后续服务的用户贡献为零。等待状态作为请求自身的一个特征，预约等待请求在 $f_A$ 中记录一次。域外随机需求作为 $\xi$ 的预测输入。

以 $s_\ell=s(t_\ell)$、$s_{\ell+H}=s(t_{\ell+H})$ 分别表示当前状态及预测终端状态。真实与预测状态采用相同特征定义和归一化参数，归一化参数由训练集确定后固定。式~\eqref{eq:state_encoding} 将完整状态映射为固定维度的网络输入，其表示能力通过第 5 节的终端价值消融实验评价。

价值函数定义为
\begin{equation}
V_\theta(s)
=d_0+\boldsymbol v_0^\top\Phi_\theta(s)
+\sum_{m=1}^{M}a_m\max\{0,\boldsymbol v_m^\top\Phi_\theta(s)+d_m\}.
\label{eq:value_network}
\end{equation}
价值函数的训练参数 $\theta$ 包含网络和状态映射中的参数。每轮优化时固定这些参数，终端库存、未完成请求、到站和截止时间特征及车辆信息则随本轮决策变化。式~\eqref{eq:terminal_vehicle_state} 中的变量乘积按附录~\ref{lin:terminal_state_section} 处理；状态编码中的有界变量乘积及 ReLU、价值网络中的 ReLU 按附录~\ref{lin:value_computation_section} 展开。由此将终端价值与业务约束共同纳入同一个优化问题。

\subsection{实际收益与状态更新}
\label{sec:reward}
每轮由优化模型返回当前路径、功率及相应服务方案，首段执行结果决定实际收益。令 $\mathcal H_\ell^{\mathrm{sw}}$ 为在 $t_\ell$ 实际获得服务的请求，$\mathcal H_\ell^{\mathrm{fail}}$ 为截止时刻位于 $\mathcal I_\ell$ 内且实际超过等待期限仍未获服务的预约请求，$N_\ell^{\mathrm{adj}}$ 为本轮实际向其发布了不同换电路径的在途用户人数，则
\begin{equation}
\begin{aligned}
r_\ell={}&E_B\sum_{r\in\mathcal H_\ell^{\mathrm{sw}}}
c_{i(r),\ell}^{\mathrm{sw}}(1-\rho_r)\\
&-\delta\sum_{i\in\mathcal S}c_{i,\ell}^{\mathrm{ch}}
\sum_{b\in\mathcal B_i}P_{i,b,\ell}^{*}
-c^{\mathrm{adj}}N_\ell^{\mathrm{adj}}
-c^{\mathrm{fail}}|\mathcal H_\ell^{\mathrm{fail}}|.
\end{aligned}
\label{eq:actual_reward}
\end{equation}
其中 $\rho_r$ 为实际退回 SOC，当前功率为执行值。实际在站请求参与首段服务。预测中的未来换电收入在实际提供服务后，计入对应时间段的收益；未满足预约请求的惩罚则在实际超过等待期限仍未提供服务时计入。跨边界等待请求在实际服务时记录一次收入；某次预约请求未获满足后，该用户的后续服务取消，惩罚只记录一次。式~\eqref{eq:actual_reward} 中的路径调整费用只按实际发布计算；尚未进入高速的用户即使在式~\eqref{eq:path_change} 中取 $\chi_{p,k}=1$，也不计入 $N_\ell^{\mathrm{adj}}$。

设 $a_\ell^*$ 汇集本轮为未进入高速的用户保存的计划，以及实际发布的在途路径、执行的功率和服务安排，以 $\varepsilon_{\ell+1}$ 表示该时间段到站记录、车辆观测及下一轮预测的更新，则
\begin{equation}
s_{\ell+1}=\mathcal F(s_\ell,a_\ell^*,\varepsilon_{\ell+1}).
\label{eq:state_transition}
\end{equation}
映射 $\mathcal F$ 按第~\ref{sec:model}~节的换电、充电、等待、路径保存及发布规则更新完整运营记录；价值评价时再按式~\eqref{eq:state_encoding} 编码。实际到站记录确定等待截止时刻；车辆完成前一站换电后，再根据实际出发时刻、该用户的预测站间行驶时间和新观测更新下一换电站的预计到站时刻。更后面的请求继续保留站点和行驶时间等信息，待其成为用户下一次换电请求时，再更新预计到站时间；车辆信息使用新的位置和 SOC 观测。因此，训练样本来自实际执行的状态和收益，并与预测终端采用相同的时间和车辆信息定义。当前时段按求解器确定的请求和电池安排执行换电，并保留相应电池变化。未进入高速的用户只更新保留的计划。

对采集的完整运营轨迹，以 $L$ 表示末端状态的索引，轨迹覆盖期间的电价、车辆观测和行驶时间信息随场景给定。剩余电池残值取零。采用未折扣回报
\begin{subequations}\label{eq:return_target}
\begin{equation}
G_\ell=\sum_{n=\ell}^{L-1}r_n,
\label{eq:return_sum}
\end{equation}
\begin{equation}
G_L=0.
\label{eq:return_terminal}
\end{equation}
\end{subequations}
式~\eqref{eq:return_sum} 只累计式~\eqref{eq:actual_reward} 中实际发生的净收益，因此 $V_\theta$ 估计的也是后续实际净收益。式~\eqref{eq:objective} 对未进入高速的用户另行计入相对日前路径的惩罚，这部分只用于优化，不进入 $r_\ell$ 或 $G_\ell$，也不作为未来实际调整费用的提前结算。

\subsection{价值函数训练}
\label{sec:training}
采用交替采样与价值拟合的训练过程。首先使用零终端价值的滚动时域优化生成完整运营轨迹。随后，以各状态对应的回报 $G_\ell$ 训练价值网络，损失函数为
\begin{equation}
\mathcal L(\theta)=\frac{1}{|\mathcal D|}
\sum_{(s_\ell,G_\ell)\in\mathcal D}
\bigl(V_\theta(s_\ell)-G_\ell\bigr)^2.
\label{eq:value_loss}
\end{equation}
其中 $\mathcal D$ 保存当前采样批次的完整状态记录与回报，训练时按当前参数 $\theta$ 重新计算 $\Phi_\theta(s_\ell)$，以同时更新状态映射和价值网络。网络更新后，将新价值函数用于下一批场景中的滚动时域优化，再采集轨迹并拟合。这样，优化模型根据当前终端价值生成决策，价值学习根据这些决策的实际结果更新评价。

\begin{table}[!htbp]
\centering
\caption{终端价值学习流程}
\label{tab:training_algorithm}
\begin{tabularx}{\linewidth}{@{}lX@{}}
\toprule 步骤 & 内容 \\ \midrule
1 & 固定训练、验证和测试场景划分，初始化用于优化的终端价值为零。\\
2 & 在一个采样批次内固定价值函数及特征提取参数，按充电功率每 5~min、两类预约用户路径每 15~min 的安排执行滚动时域优化，记录完整运营轨迹。\\
3 & 由式~\eqref{eq:actual_reward} 计算各段实际收益，按式~\eqref{eq:return_target} 从场景结束向前累计回报。\\
4 & 最小化式~\eqref{eq:value_loss}，更新状态映射和价值网络；归一化统计在训练阶段确定并固定。\\
5 & 用更新后的价值函数重复采样和拟合，按验证场景的实际净收益选择网络。\\
6 & 固定选定网络，在独立测试场景上评价运营表现和求解时间。\\
\bottomrule
\end{tabularx}
\end{table}

该过程可视为由滚动时域优化生成策略、由回报拟合进行评价的近似策略迭代。每次求解期间价值函数保持固定，参数更新发生在采样批次之间。训练场景覆盖不同需求强度、SOC 分布和等待情况，网络结构、批次规模和训练预算在实验设置中记录。价值近似与滚动优化之间的相互作用通过验证场景的实际运营表现评价。


% !TeX root = ../main.tex
\section{数值实验}
\label{exp:section}

数值实验从换电路径与充电功率的联合优化、RL 终端价值和运行条件三个方面评价所提方法。本次完成第~\ref{exp:terminal}~节的零终端价值实验；RL 训练、RL 对照及完美信息结果尚未开展。各预测域使用相同的道路网络、七个测试日、预约与随机请求、误差、日前路径和初始库存，需求预测仅使用决策时刻可获得的信息。

道路包含 11 个换电站、6 个出入口和 26 个 O--D 对。节点名称与坐标来自 \path{poi_node_6e.csv}，站点与出入口里程分别来自 \path{distance_stations.csv}、\path{distance_6e_entrance.csv}；站点电池配置、电价和服务费分别来自 \path{11s_configuration_info.csv}、\path{11s_elec_price_info.csv}、\path{11s_service_fee_info.csv}。O--D 与人口数据来自 \path{subpaths_6e.csv} 和 \path{population_6e.csv}。预约 O--D 抽样权重为起终点人口乘积除以距离平方，入场小时权重来自 \path{flowVolume.csv}，小时内入场时刻均匀抽样。

根据 \texttt{demand\_info.csv} 中站点、小时和星期对应的真实需求均值，独立进行泊松抽样，生成连续十周、共 70 天的\textbf{基于真实均值的泊松合成场景}。从每种星期对应的十天中随机选一天，周一至周日的测试日编号依次为 1、37、3、4、19、13、14；其余 63 天作为后续训练和验证候选池，尚未进一步划分。每个测试日均生成 200 名预约用户和 200 个随机请求。随机请求以该日泊松样本的站点与小时联合频数为权重进行固定总量的多项分配，再在小时内均匀抽样到达时刻；零总量样本报错，不更换测试日。随机需求预测使用原文件中对应星期的均值，归一化至每日期望 200，并按累计期望量生成确定性预测请求，不读取测试日未来实际请求。数据来源、随机流及场景哈希保存于 \texttt{data/final\_real\_data\_v1/manifest.json}；各随机流均由基础种子 1 按日期和用途派生。

计算环境见表~\ref{exp:environment}，主要参数见表~\ref{exp:parameters}。本章实验取 $\delta=15$~min、$K=2$，覆盖模型说明中默认的 5~min 和 $K=3$：充电功率每段更新，换电路径每 30~min 允许更新。$H=4,5,6,7,8$ 分别对应 60、75、90、105、120~min，$H$ 本身即时间段数。单次求解时限为 120~s，目标相对 gap 为 $10^{-4}$（0.01\%），每个日任务最多使用 16 线程，各任务串行运行。达到时限但存在可行解时执行该解，并保留实际 gap；没有可行解时停止任务并保存诊断。

\begin{table}[!htbp]
  \centering
  \caption{计算环境}
  \label{exp:environment}
  \small
  \begin{tabularx}{0.88\linewidth}{@{}lX@{}}
    \toprule
    项目 & 配置 \\
    \midrule
    CPU & Intel Core i7-12650H，10 核、16 逻辑线程 \\
    系统可见内存 & 15.7~GiB \\
    操作系统 & Windows，系统构建号 10.0.26200 \\
    Python 版本 & 3.10.19 \\
    COPT 版本 & 8.0.6 \\
    GPU / PyTorch & 本次零终端价值实验未使用 \\
    \bottomrule
  \end{tabularx}
\end{table}

\begin{table}[!htbp]
  \centering
  \caption{主要实验参数}
  \label{exp:parameters}
  \small
  \setlength{\tabcolsep}{5pt}
  \begin{tabularx}{\linewidth}{@{}Xll@{}}
    \toprule
    参数 & 符号 & 设置 \\
    \midrule
    站点数量 & $|\mathcal S|$ & 11 \\
    出入口与 O--D 数量 & --- & 6，26 \\
    电池总数与初始状态 & --- & 250 块，全部满电 \\
    电池容量 & $E_B$ & 100~kWh \\
    各站充电效率 & $\eta_i$ & 95\% \\
    设备功率上限 & $\overline P_{i,b}^{\mathrm{ch}}$ & 60~kW \\
    站点总充电功率上限 & $\overline P_i^{\mathrm{sta}}$ & 960~kW \\
    用户最长等待时间 & --- & 30~min \\
    路径调整与预约失败成本 & $c^{\mathrm{adj}},c^{\mathrm{fail}}$ & 10 元/次，1000 元/人 \\
    充电更新周期 & $\delta$ & 15~min \\
    路径更新周期 & $K\delta$ & 30~min \\
    预测域长度 & $H$ & 4、5、6、7、8 段 \\
    需求期长度 & $N$ & 96 段，24~h \\
    日需求量：预约 / 随机 & --- & 200 / 200 \\
    车速、满电续航 & --- & 75~km/h，300~km \\
    出口 SOC 下限、间距剪枝阈值 & --- & 0.1，100~km \\
    本次终端价值权重 & $\beta$ & 0 \\
    候选池 / 测试日数量 & --- & 63 / 7 \\
    单次求解时限、目标相对 gap & --- & 120~s，0.01\% \\
    基础随机种子 & --- & 1 \\
    \bottomrule
  \end{tabularx}
\end{table}

按 $s_1$ 至 $s_{11}$ 顺序，各站电池数为 $[22,22,22,18,20,25,23,20,28,28,22]$；变量、执行和统计均按实际槽数处理。预约实际入口 SOC 在 $(0.5,1]$ 上均匀生成，上报 SOC 限于 $[0.5,1]$；上报时间误差上限为 10~min，SOC 误差上限为 0.05，在各自可行区间内均匀生成。物理路段行驶时间乘数独立取自 $U[0.9,1.1]$；随机退回 SOC 取自 $U[0.1,0.3]$，预测值取 0.2。为保证入口误差下的首段可达性，日前及未入场用户的候选网络使用 $\max(0.5,s^{\mathrm{report}}-0.05)$ 作为可达性计算的入口 SOC；预测退回电量使用报告值，实际执行使用真实值。100~km 沿用第~\ref{sec:network}~节的剪枝规则，保留必要可行路径和既有计划，并非对所有弧施加硬下限。

24:00 后停止新增预约与随机需求，随机预测同样不再生成新请求；预测窗口保持指定 $H$，已有用户继续处理至全部完成服务或失败。预约完成按其剩余换电安排为空且可满足出口 SOC 要求判定，最后换电后的无服务行驶不再延长仿真。价格按 24 小时循环，跨日收入和费用均归入原需求日，不计库存残值、不强制补满。按最长行程、行驶误差和逐站最大等待推导的安全上限为 186 段（46.5~h）；到达该上限仍有待处理用户时标记异常。需求期、实际执行期和跨日处理期分别记录。

运营净收益为实际服务收入减去充电、路径调整和预约失败成本。换电收入按实际补充电量乘以当时电价与服务费之和计算，充电成本按电网输入电量另行扣除。预约失败率为失败用户数除以 200，同一用户仅罚一次；随机服务率为服务请求数除以 200。平均等待时间先在每个日场景内对全部已服务换电请求计算，再对七天取平均。路径调整次数仅记录向在途用户实际发布的不同换电站序列，未入场计划变化不计实际费用。

所有主指标均报告固定七个测试日的算术均值；每个 $H$ 每天仅运行一次，共 35 个日任务，属于同一基础种子下的跨日比较，不是多种子重复实验。表~\ref{exp:value_results} 的求解均值及 P95 先在每日需求期的 96 轮内计算，再对七天取平均；含跨日阶段的全部轮次统计另行报告。P95 为日内第 95 百分位数，七天不按轮次数加权。样本分母为零时记为“不适用”。



\FloatBarrier
\subsection{The role of RL terminal value function}
\label{exp:terminal}

本节拟比较 $\beta=0$ 的无终端价值设置与 $\beta=1$ 的 RL 终端价值设置，预测域取 $H=4,\ldots,8$。当前已完成零终端价值的 35 个日场景；其结果刻画短预测域下的运营表现和计算代价，尚不能单独证明 RL 终端价值的收益。

终端价值网络的训练设置保留在表~\ref{exp:training}，待实际训练后填写。表~\ref{exp:value_results} 按 $H$ 列出已完成的零终端价值结果，RL 及完美信息对照保持待填。

\begin{table}[!htbp]
  \centering
  \caption{终端价值网络的训练设置}
  \label{exp:training}
  \small
  \begin{tabularx}{\linewidth}{@{}Xc@{}}
    \toprule
    项目 & 设置 \\
    \midrule
    输入维度 & 待填 \\
    隐藏单元数 & 待填 \\
    参数总数 & 待填 \\
    训练场景数量 & 待填 \\
    采样批次数与每批轨迹数 & 待填 \\
    每批参数更新次数 & 待填 \\
    回归批量大小 & 待填 \\
    优化器与学习率 & 待填 \\
    独立训练次数 & 待填 \\
    验证频率与网络选择规则 & 待填 \\
    实际训练耗时 & 待填 \\
    \bottomrule
  \end{tabularx}
\end{table}

\begin{table}[!htbp]
  \centering
  \caption{终端价值实验的运营结果（固定七日算术均值）}
  \label{exp:value_results}
  \small
  \setlength{\tabcolsep}{3pt}
  \begin{tabularx}{\linewidth}{@{}Xcccccc@{}}
    \toprule
    方法 & \shortstack{净收益\\元} & \shortstack{预约\\失败率/\%} & \shortstack{随机用户\\服务率/\%} & \shortstack{等待时间\\min} & \shortstack{调整次数\\次} & \shortstack{求解时间\\均值/P95，s} \\
    \midrule
    零终端 $H=4$ & $-96706.58$ & 60.57 & 62.36 & 11.94 & 45.71 & 0.11/0.23 \\
    零终端 $H=5$ & $-83734.73$ & 55.00 & 64.07 & 11.53 & 54.00 & 0.18/0.42 \\
    零终端 $H=6$ & $-49887.07$ & 40.07 & 70.29 & 12.10 & 76.00 & 1.13/2.06 \\
    零终端 $H=7$ & 37180.19 & 0.21 & 86.93 & 12.87 & 141.14 & 23.09/120.01 \\
    零终端 $H=8$ & 37639.56 & 0.43 & 89.57 & 12.55 & 198.71 & 50.15/120.19 \\
    \midrule
    RL $H=4$ & 待填 & 待填 & 待填 & 待填 & 待填 & 待填 \\
    RL $H=5$ & 待填 & 待填 & 待填 & 待填 & 待填 & 待填 \\
    RL $H=6$ & 待填 & 待填 & 待填 & 待填 & 待填 & 待填 \\
    RL $H=7$ & 待填 & 待填 & 待填 & 待填 & 待填 & 待填 \\
    RL $H=8$ & 待填 & 待填 & 待填 & 待填 & 待填 & 待填 \\
    完美信息上界 & 待填 & 待填 & 待填 & 待填 & 待填 & 待填 \\
    \bottomrule
  \end{tabularx}
\end{table}

\begin{table}[!htbp]
  \centering
  \caption{零终端价值的求解与跨日处理诊断}
  \label{exp:zero_diagnostics}
  \small
  \setlength{\tabcolsep}{3pt}
  \begin{tabularx}{\linewidth}{@{}cXXXXXX@{}}
    \toprule
    $H$ & \shortstack{限时轮次\\比例/\%} & \shortstack{gap 达标\\比例/\%} & \shortstack{平均实际\\gap/\%} & \shortstack{跨日时长\\h} & \shortstack{终态库存\\kWh} & \shortstack{终态满电\\电池数} \\
    \midrule
    4 & 0.00 & 100.00 & 0.000052 & 3.11 & 7054.81 & 12.29 \\
    5 & 0.00 & 100.00 & 0.000163 & 5.00 & 6609.73 & 12.00 \\
    6 & 0.24 & 99.76 & 0.000899 & 5.82 & 5512.72 & 6.00 \\
    7 & 10.19 & 89.81 & 0.163012 & 11.18 & 4864.18 & 0.71 \\
    8 & 26.62 & 73.38 & 0.969465 & 11.57 & 4735.21 & 1.14 \\
    \bottomrule
  \end{tabularx}
  \par\smallskip
  \parbox{\linewidth}{\footnotesize 注：各项先按日统计，再对七日等权平均。此表的求解比例及 gap 使用包含跨日处理的全部轮次；gap 达标指实际相对 gap 不超过 0.01\%。终态满电数量为日终计数的均值，故可为小数。}
\end{table}

当 $H$ 从 4 增至 7 时，七日平均净收益由 $-96706.58$ 元增至 37180.19 元，预约失败率由 60.57\% 降至 0.21\%，随机服务率由 62.36\% 增至 86.93\%。$H=4$ 时每日平均充电费用仅为 662.48 元，而预约失败费用为 121142.86 元。以随机请求的预测退回 SOC 0.2 为例，在 60~kW、95\% 效率下将 100~kWh 电池充满需要约 84.21~min，即至少六个完整充电时段，之后的边界才能再次服务。较短窗口难以同时看到这笔充电成本及其后续服务收益；$H=7$ 可覆盖六段充电后的服务边界。这为收益在 $H=6$ 与 $H=7$ 之间明显变化提供了物理解释，但不构成对所有请求或场景的阈值保证。

$H=8$ 相比 $H=7$ 的七日平均净收益增加 459.37 元（约 1.24\%），随机服务率提高 2.64 个百分点，需求期单轮平均求解时间则由 23.09~s 增至 50.15~s。预约失败总人数由七天合计 3 人增至 6 人，因此并非所有服务指标均随 $H$ 单调改善。需求期每轮二元变量数的七日均值由 $H=4$ 的约 2489 个增至 $H=8$ 的约 3933 个；限时求解比例也明显上升。含跨日阶段时，$H=4,\ldots,8$ 的单轮求解均值/P95 七日均值依次为 0.10/0.22、0.15/0.40、0.90/1.58、15.79/120.00、34.08/120.15~s，低于仅需求期统计的相应均值。

表~\ref{exp:zero_diagnostics} 表明，0.01\% 是求解目标，不能表述为全部轮次均达到的精度。$H=7$ 和 $H=8$ 全部日场景中的最大实际 gap 分别为 6.02\% 和 16.97\%。这些有限时长可行解对应的运营轨迹均完成核验，结果比较仍受求解预算影响，不能将差异全部归因于预测域长度。求解器退出和计时开销使少数记录略高于 120~s。

35 条轨迹共包含 4387 个执行时段。独立重建账务及指标后，每日预约完成数加失败数均为 200，随机服务数加超时数均为 200，最终无未处理用户。满电交付、电池 SOC 与功率、预约优先及同类先到先服务、后站到达对前站服务的依赖、等待截止边界及实际发布计费均通过核验。跨日处理最长为 14.5~h，低于推导的安全上限；道路最长 O--D 为 1140~km，长行程与晚间入场可产生较长跨日服务。当前结论仅适用于初始全满且无库存残值结算的独立需求日，不能直接解释为连续多日稳态收益。

\begin{figure}[!htbp]
  \centering
  \fbox{\begin{minipage}[c][32mm][c]{0.9\linewidth}
    \centering
    待填：价值网络训练过程及终端价值消融结果\\[4pt]
    横轴与单位、RL 对照曲线及按测试日统计的差异待填。
  \end{minipage}}
  \caption{终端价值学习过程及运营表现（待填）}
  \label{exp:value_figure}
\end{figure}

\todoresult{待填：完成 RL 训练及配对测试后，比较相同预测域下的各类终端价值；完美信息结果另行计算，不从零终端价值结果推断。}


\FloatBarrier
\subsection{The value of route and charging power optimization}
\label{exp:joint}

为考察两项决策的作用，设置表~\ref{exp:joint_methods}所列四组方法。各组终端价值均取零，使用相同预测域长度、需求输入和服务规则。服务时间与满电电池匹配均由优化模型确定。采用日前路径的组别在运营过程中沿用日前计算的换电站序列；优化换电路径的组别按规定的更新周期调整在途预约用户的换电安排。

\begin{table}[!htbp]
  \centering
  \caption{换电路径与充电功率联合优化的对照设置}
  \label{exp:joint_methods}
  \small
  \begin{tabularx}{\linewidth}{@{}Xlll@{}}
    \toprule
    方法 & 换电路径 & 充电功率 & 终端价值 \\
    \midrule
    基准方法 & 日前路径 & 基准充电规则 & 0 \\
    仅优化换电路径 & 在线优化 & 基准充电规则 & 0 \\
    仅优化充电功率 & 日前路径 & 在线优化 & 0 \\
    联合优化 & 在线优化 & 在线优化 & 0 \\
    完美信息上界 & 在线优化 & 在线优化 & 0 \\
    \bottomrule
  \end{tabularx}
\end{table}

基准充电规则将站点总充电功率上限按设备数量等分，实际功率同时满足设备上限和电池可充电量限制：
\begin{equation}
 P_{i,b,n}^{\mathrm{base}}=
 \min\left\{
 \overline P_{i,b}^{\mathrm{ch}},\;
 \frac{\overline P_i^{\mathrm{sta}}}{|\mathcal B_i|},\;
 \frac{E_B(1-S_{i,b,n}^{+})}{\eta_i\delta}
 \right\}.
 \label{eq:baseline_power}
\end{equation}
该规则使用本时刻完成换电后的 SOC。每台设备保留等分得到的功率额度，满电电池对应设备的功率为零。采用基准规则的组别在预测域内施加 $P_{i,b,n}=(1-d_n^{\mathrm{end}})P_{i,b,n}^{\mathrm{base}}$，使各运营时间段使用基准功率，结束后的时间段功率为零。其分段线性表达见附录。式中 $\delta$ 以小时计，因而第三项的单位为 kW。

表~\ref{exp:joint_results}记录运营指标，表~\ref{exp:joint_costs}记录收入与成本分解。通过基准方法与两组单项优化的比较，考察路径调整和充电安排各自的作用；通过两组单项优化与联合优化的比较，考察同时优化两项决策的运营表现。各项比较使用相同测试场景的配对结果。

\begin{table}[!htbp]
  \centering
  \caption{联合优化消融实验的运营结果}
  \label{exp:joint_results}
  \small
  \setlength{\tabcolsep}{3pt}
  \begin{tabularx}{\linewidth}{@{}Xcccccc@{}}
    \toprule
    方法 & \shortstack{净收益\\元} & \shortstack{预约\\失败率/\%} & \shortstack{随机用户\\服务率/\%} & \shortstack{等待时间\\min} & \shortstack{调整次数\\次} & \shortstack{求解时间\\均值/P95，s} \\
    \midrule
    基准方法 & 待填 & 待填 & 待填 & 待填 & 待填 & 待填 \\
    仅优化换电路径 & 待填 & 待填 & 待填 & 待填 & 待填 & 待填 \\
    仅优化充电功率 & 待填 & 待填 & 待填 & 待填 & 待填 & 待填 \\
    联合优化 & 待填 & 待填 & 待填 & 待填 & 待填 & 待填 \\
    完美信息上界 & 待填 & 待填 & 待填 & 待填 & 待填 & 待填 \\
    \bottomrule
  \end{tabularx}
\end{table}

\begin{table}[!htbp]
  \centering
  \caption{联合优化消融实验的收入与成本分解（元）}
  \label{exp:joint_costs}
  \small
  \begin{tabularx}{\linewidth}{@{}Xcccc@{}}
    \toprule
    方法 & 服务收入 & 充电成本 & 路径调整成本 & 预约失败成本 \\
    \midrule
    基准方法 & 待填 & 待填 & 待填 & 待填 \\
    仅优化换电路径 & 待填 & 待填 & 待填 & 待填 \\
    仅优化充电功率 & 待填 & 待填 & 待填 & 待填 \\
    联合优化 & 待填 & 待填 & 待填 & 待填 \\
    完美信息上界 & 待填 & 待填 & 待填 & 待填 \\
    \bottomrule
  \end{tabularx}
\end{table}

\todoresult{待填：根据运营指标和成本分解，分析换电路径优化、充电功率优化及二者联合优化的作用，并给出重复实验的统计结果。}


\FloatBarrier
\subsection{敏感性分析}
\label{exp:sensitivity}

敏感性分析考察需求强度、预约和随机人数的比例、电池库存配置、用户旅行时间的波动性。表~\ref{exp:sensitivity_settings}记录各因素的基准值和变化范围。

\begin{table}[!htbp]
  \centering
  \caption{敏感性分析的参数设置}
  \label{exp:sensitivity_settings}
  \small
  \begin{tabularx}{\linewidth}{@{}Xlll@{}}
    \toprule
    因素 & 基准值 & 变化范围 & 取值数量 \\
    \midrule
    需求强度 & 待填 & 待填 & 待填 \\
    预约与随机用户比例 & 待填 & 待填 & 待填 \\
    各站电池数量 $|\mathcal B_i|$ & 待填 & 待填 & 待填 \\
    用户旅行时间波动性 & 待填 & 待填 & 待填 \\
    \bottomrule
  \end{tabularx}
\end{table}



\begin{table}[!htbp]
  \centering
  \caption{敏感性分析结果（各参数取值分别填列）}
  \label{exp:sensitivity_results}
  \small
  \setlength{\tabcolsep}{3pt}
  \begin{tabularx}{\linewidth}{@{}Xcccccc@{}}
    \toprule
    因素、取值及方法 & \shortstack{净收益\\元} & \shortstack{预约\\失败率/\%} & \shortstack{随机用户\\服务率/\%} & \shortstack{等待时间\\min} & \shortstack{调整次数\\次} & \shortstack{求解时间\\均值/P95，s} \\
    \midrule
    需求强度：待填 & 待填 & 待填 & 待填 & 待填 & 待填 & 待填 \\
    预约与随机用户比例：待填 & 待填 & 待填 & 待填 & 待填 & 待填 & 待填 \\
    电池库存配置：待填 & 待填 & 待填 & 待填 & 待填 & 待填 & 待填 \\
    用户旅行时间波动性：待填 & 待填 & 待填 & 待填 & 待填 & 待填 & 待填 \\
    \bottomrule
  \end{tabularx}
\end{table}

\begin{figure}[!htbp]
  \centering
  \fbox{\begin{minipage}[c][36mm][c]{0.9\linewidth}
    \centering
    待填：预测域长度、需求强度、电池库存和站点总充电功率上限的敏感性结果\\[4pt]
    各参数取值、运营指标曲线及求解时间曲线待填。
  \end{minipage}}
  \caption{不同运行条件下的运营表现与求解时间（待填）}
  \label{exp:sensitivity_figure}
\end{figure}

\todoresult{待填：分别分析预测域长度、需求强度、电池数量和站点总充电功率上限变化时的收益、服务表现及求解时间，并说明各项结论所对应的参数范围。}


% !TeX root = ../main.tex
\section{结论}

本文围绕高速公路换电站网络中的预约与随机用户，建立换电路径与充电功率的滚动时域优化方法。候选网络根据 SOC 和道路条件描述车辆可达的换电安排，日前路径提供在线调整的基准。优化模型将路径选择、充电、电池状态及用户服务条件联系起来，以不同周期更新换电路径和充电功率，并按实际服务、充电、路径改变和预约失败记录运营收益。

强化学习用于估计预测域的终端价值，将终端库存与预测域外未完成服务共同纳入评价。模型保留电池 SOC 分布、等待信息及预约后续换电需求，并通过分段线性网络表达其价值。价值函数由完整运营轨迹中的实际回报训练，网络参数在每轮优化中固定。数值实验按联合优化消融、RL 终端价值消融和敏感性分析组织，具体结果与定量结论将在实验完成后填写。

本文以每轮 ETA 预测描述到站窗口，以瞬时换电和给定充电效率表达站内运行。状态编码和价值函数具有有限维度，其效果需要结合实际运营指标和求解时间评价。后续研究可进一步结合交通过程、换电作业时间及电池特性，分析更丰富的运行条件下路径与充电的协调方法。



\appendix
% !TeX root = ../main.tex
\section{正文各部分的线性化方法}
\label{lin:appendix}

本附录按正文对应部分说明线性化方法，各小节开头给出相应公式编号。流平衡、换电次数、到站时刻与等待期限、满电交付和充电等关系已按线性形式给出，可直接使用。每轮求解时，网络参数和各表达式的有效上下界均为已知常数，新增变量仅用于表达相应关系。

\subsection{路径调整约束线性化}
\label{lin:path_adjustment_section}
本节对应正文的“路径调整约束”，给出式~\eqref{eq:path_freeze} 和式~\eqref{eq:path_change} 的等价线性表达。

对各 $p\in\mathcal P$ 和 $i\in\mathcal S$，式~\eqref{eq:path_freeze} 中的绝对值不等式可直接写为
\begin{equation}
\begin{aligned}
-g_\ell\le x_i^{p,k}-\widehat{x}_i^{p,k}&\le g_\ell,
&&k\in\mathcal K^{\mathrm{fut},p},\\
-g_\ell\le x_i^{p,k}-x_i^{p,k,\mathrm{pub}}&\le g_\ell,
&&k\in\mathcal K^{\mathrm{enr},p},
\end{aligned}
\label{lin:path_freeze}
\end{equation}
其中，$\widehat{x}_i^{p,k}$ 表示尚未进入高速公路的用户在上一轮保留的未发布计划，初始取日前安排 $\bar x_i^{p,k}$；$x_i^{p,k,\mathrm{pub}}$ 表示在途用户最近一次已发布路径中尚未完成的换电安排。两者在本轮求解时均已知。比较范围与正文一致，各条路径均去除已经完成的换电和当前正在等待的换电。$g_\ell=0$ 时，式~\eqref{lin:path_freeze} 要求保持这些安排；$g_\ell=1$ 时，允许调整换电站点。

计算式~\eqref{eq:path_change} 中的路径改变标记时，尚未进入高速公路的用户与日前路径比较，在途用户与最近一次已发布的剩余路径比较。为简化本附录的表达，记
\begin{equation}
x_i^{p,k,\mathrm{ref}}=
\begin{cases}
\bar x_i^{p,k}, & k\in\mathcal K^{\mathrm{fut},p},\\
x_i^{p,k,\mathrm{pub}}, & k\in\mathcal K^{\mathrm{enr},p},
\end{cases}
\label{lin:path_reference}
\end{equation}
对上述两类用户的每个站点引入连续变量 $D_i^{p,k}$，用以下约束表示 $D_i^{p,k}=|x_i^{p,k}-x_i^{p,k,\mathrm{ref}}|$：
\begin{equation}
\begin{aligned}
D_i^{p,k}&\ge x_i^{p,k}-x_i^{p,k,\mathrm{ref}},&
D_i^{p,k}&\ge x_i^{p,k,\mathrm{ref}}-x_i^{p,k},\\
D_i^{p,k}&\le x_i^{p,k}+x_i^{p,k,\mathrm{ref}},&
D_i^{p,k}&\le 2-x_i^{p,k}-x_i^{p,k,\mathrm{ref}},\\
0&\le D_i^{p,k}\le1.
\end{aligned}
\label{lin:path_difference}
\end{equation}
由于 $x_i^{p,k}$ 和 $x_i^{p,k,\mathrm{ref}}$ 均取 $0$ 或 $1$，式~\eqref{lin:path_difference} 在两者相同时要求 $D_i^{p,k}=0$，在两者不同时要求 $D_i^{p,k}=1$。对各 $p\in\mathcal P$ 和 $k\in\mathcal K^{\mathrm{fut},p}\cup\mathcal K^{\mathrm{enr},p}$，再加入
\begin{equation}
\begin{aligned}
\chi_{p,k}&\ge D_i^{p,k}\quad(i\in\mathcal S),\\
\chi_{p,k}&\le\sum_{i\in\mathcal S}D_i^{p,k},
\qquad \chi_{p,k}\in\{0,1\},
\end{aligned}
\label{lin:path_change}
\end{equation}
即可得到式~\eqref{eq:path_change} 的等价线性约束。式~\eqref{lin:path_change} 要求：只要有一个站点的换电安排与比较路径不同，$\chi_{p,k}$ 就取 $1$；所有站点的安排均相同时，$\chi_{p,k}$ 取 $0$。因此，即使目标函数中的路径调整费用系数为零，这些约束仍能确定正确的路径改变标记。

\subsection{等待和服务顺序约束线性化}
\label{lin:waiting_priority_section}
本节对应正文的“等待和服务顺序”，处理式~\eqref{eq:queue_status} 中的变量乘积和式~\eqref{eq:priority} 中的服务先后条件。

由换电次数约束，请求在任意时刻之前的换电总次数只能为 $0$ 或 $1$。因此，对所有 $r\in\R$ 和 $n\in\mathcal T_\ell\cup\{\ell+H\}$，式~\eqref{eq:queue_status} 等价于
\begin{equation}
\begin{aligned}
0\le Q_{r,n}&\le a_r,\\
Q_{r,n}&\le e_{r,n},\\
Q_{r,n}&\le1-
\sum_{\substack{m\in\mathcal T_\ell\\m<n}}
\sum_{b\in\mathcal B_{i(r)}}\alpha_{b,r,m},\\
Q_{r,n}&\ge a_r+e_{r,n}-1-
\sum_{\substack{m\in\mathcal T_\ell\\m<n}}
\sum_{b\in\mathcal B_{i(r)}}\alpha_{b,r,m}.
\end{aligned}
\label{lin:queue_status}
\end{equation}
三个条件全部满足时，上述约束要求 $Q_{r,n}=1$；任一条件不满足时，要求 $Q_{r,n}=0$。当前时刻 $n=\ell$ 的求和为空，预测终点的求和则包含本轮全部换电。

对于同站的预约请求和随机请求，正文式~\eqref{eq:priority} 已是线性约束，可直接使用。

对于同站同类的两个不同请求 $q,r$，令 $M_i^{\mathrm{arr}}$ 为本站所有候选到站时刻的最大值与最小值之差。候选值包括给定到站时刻、前站在各 $t_m$ 换电后得到的 $t_m+\tau_{j,i}^{p,k}$，以及时间尚未确定时的零值。该常数在本轮求解前计算。对所有同站同类的有序请求对和 $n\in\mathcal T_\ell$，施加
\begin{equation}
t_r^{\mathrm{arr}}-t_q^{\mathrm{arr}}
\le M_i^{\mathrm{arr}}\left(
2-Q_{q,n}
-\sum_{b\in\mathcal B_i}\alpha_{b,r,n}
+\sum_{b\in\mathcal B_i}\alpha_{b,q,n}
\right).
\label{lin:arrival_priority}
\end{equation}
只有请求 $r$ 本次获得服务、请求 $q$ 正在等待但本次未获服务时，右侧括号内才为 $0$，此时要求 $r$ 不晚于 $q$ 到站。其他情况下，右侧至少为 $M_i^{\mathrm{arr}}$，约束自动满足。因此，式~\eqref{lin:arrival_priority} 与正文的先到先服务条件等价；同时到站时不规定先后，尚未到站或到站时间尚未确定的请求也不会阻止其他请求获得服务。

\subsection{未满足预约请求判定约束线性化}
\label{lin:unmet_requests_section}
本节对应正文的“未满足预约请求的判定”。式~\eqref{eq:deadline_passed} 已是线性表达；需要转换的是式~\eqref{eq:failure_exact} 中的变量乘积。

对所有 $r\in\R^A$，式~\eqref{eq:failure_exact} 等价于
\begin{equation}
\begin{aligned}
0\le f_r&\le a_r,\\
f_r&\le h_r,\\
f_r&\le1-\sum_{n\in\mathcal T_\ell}
\sum_{b\in\mathcal B_{i(r)}}\alpha_{b,r,n},\\
f_r&\ge a_r+h_r-1-\sum_{n\in\mathcal T_\ell}
\sum_{b\in\mathcal B_{i(r)}}\alpha_{b,r,n}.
\end{aligned}
\label{lin:unmet_reservation}
\end{equation}
这些约束要求：只有请求需要服务、截止早于预测终点，且本轮未完成换电时，$f_r$ 才取 $1$。

\subsection{预测终端状态的线性化}
\label{lin:terminal_state_section}
本节对应第~\ref{sec:value}~节的“预测终端状态”，处理式~\eqref{eq:terminal_pending}、式~\eqref{eq:terminal_random} 和式~\eqref{eq:terminal_vehicle_state} 中的变量乘积。

对预约请求 $r\in\R_{p,k}^A$，式~\eqref{eq:terminal_pending} 等价于
\begin{equation}
\begin{aligned}
0\le w_r&\le a_r,\\
w_r&\le1-\sum_{n\in\mathcal T_\ell}
\sum_{b\in\mathcal B_{i(r)}}\alpha_{b,r,n},\\
w_r&\le1-\sum_{q\in\R_{p,k}^A}f_q,\\
w_r&\ge a_r-\sum_{n\in\mathcal T_\ell}
\sum_{b\in\mathcal B_{i(r)}}\alpha_{b,r,n}
-\sum_{q\in\R_{p,k}^A}f_q.
\end{aligned}
\label{lin:terminal_pending}
\end{equation}
对随机请求 $r\in\R^R$，式~\eqref{eq:terminal_random} 等价于
\begin{equation}
\begin{aligned}
0\le w_r&\le1-h_r,\\
w_r&\le1-\sum_{n\in\mathcal T_\ell}
\sum_{b\in\mathcal B_{i(r)}}\alpha_{b,r,n},\\
w_r&\ge1-h_r-\sum_{n\in\mathcal T_\ell}
\sum_{b\in\mathcal B_{i(r)}}\alpha_{b,r,n}.
\end{aligned}
\label{lin:terminal_random}
\end{equation}
两组约束分别确定预测终点仍需服务的预约请求和随机请求。

式~\eqref{eq:terminal_vehicle_state} 中，候选车辆信息是已知常数。对其中出现的乘积 $w_r\alpha_{b,q,m}$，引入辅助量 $u_{b,q,r,m}$，并施加
\begin{equation}
\begin{aligned}
0\le u_{b,q,r,m}&\le w_r,\\
u_{b,q,r,m}&\le\alpha_{b,q,m},\\
u_{b,q,r,m}&\ge w_r+\alpha_{b,q,m}-1.
\end{aligned}
\label{lin:terminal_vehicle_product}
\end{equation}
该式适用于 $r\in\R^A$、$q\in\mathcal P_r$、$m\in\mathcal T_\ell$ 和 $b\in\mathcal B_{i(q)}$。将相应乘积替换为 $u_{b,q,r,m}$ 后，终端车辆信息即为辅助量与已知候选值的线性组合。终端电池电量、等待标记、到站时刻以及路径记录继续使用正文已有关系。

\subsection{状态编码与价值计算的线性化}
\label{lin:value_computation_section}
本节对应第~\ref{sec:value}~节的“状态编码与价值计算”，处理式~\eqref{eq:state_encoding} 中请求标记与特征的乘积，以及式~\eqref{eq:value_network} 和特征网络中的分段线性函数。

设 $x\in[L,U]$、$b\in\{0,1\}$，引入连续变量 $w=bx$，其等价表达为
\begin{equation}
\begin{aligned}
Lb\le w&\le Ub,\\
x-U(1-b)\le w&\le x-L(1-b).
\end{aligned}
\label{lin:product}
\end{equation}
当 $b=0$ 时得到 $w=0$，当 $b=1$ 时得到 $w=x$。该表达适用于含负值的区间，可用于请求状态标记与连续特征的乘积。

对 $h=\max\{0,x\}$，若 $L<0<U$，引入二元变量 $\xi$，有
\begin{equation}
\begin{aligned}
h&\ge 0, & h&\ge x,\\
h&\le U\xi, & h&\le x-L(1-\xi),
\qquad \xi\in\{0,1\}.
\end{aligned}
\label{lin:relu}
\end{equation}
当 $U\le0$ 时直接置 $h=0$；当 $L\ge0$ 时直接置 $h=x$。对式~\eqref{eq:value_network} 中每个仿射输入 $x_m=\boldsymbol v_m^\top\Phi_\theta(s)+d_m$ 分别使用上述约束，并将对应的输出变量 $h_m$ 代入
\begin{equation}
V_\theta(s)=d_0+\boldsymbol v_0^\top\Phi_\theta(s)+\sum_m a_m h_m.
\label{lin:network}
\end{equation}
上述等式和约束共同确定网络值，适用于任意符号的输出权重 $a_m$。状态编码中的分段线性特征按相同方法展开；其输出作为后续网络的输入，从而将特征提取与价值计算一并纳入优化模型。

上述乘积和 ReLU 约束所需的上下界按以下方式确定。
SOC 的边界取 $[0,1]$，请求状态标记取 $[0,1]$，功率使用正文中的设备和站级上限。请求数量由当前有限请求集合限定；到站时刻及与等待期限有关的时间特征根据本轮已知时刻及有效等待窗口确定范围。对于含状态标记的特征，先求连续部分的边界，再使用式~\eqref{lin:product}。若某层输入向量为 $\boldsymbol e$，$x=\boldsymbol v^\top\boldsymbol e+d$ 且 $e_j\in[\underline e_j,\overline e_j]$，令 $v_j^+=\max\{v_j,0\}$、$v_j^-=\min\{v_j,0\}$，则可取
\begin{equation}
\begin{aligned}
L&=d+\sum_j\left(v_j^+\underline e_j+v_j^-\overline e_j\right),\\
U&=d+\sum_j\left(v_j^+\overline e_j+v_j^-\underline e_j\right).
\end{aligned}
\label{lin:interval}
\end{equation}
ReLU 输出的区间为 $[\max\{0,L\},\max\{0,U\}]$。依次传递各层区间，可得到隐藏单元和最终输出的有限上下界，并据此固定本轮线性化常数。若求解前能够由业务约束得到更紧的边界，则使用相应边界以缩小连续松弛范围。

\subsection{基准充电规则线性化}
\label{lin:baseline_charging_section}
对式~\eqref{eq:baseline_power}，记已知常数与仿射表达式为
\begin{equation}
\begin{aligned}
\bar p_{i,b}
&=\min\left\{\overline P_{i,b}^{\mathrm{ch}},
\frac{\overline P_i^{\mathrm{sta}}}{|\mathcal B_i|}\right\},\\
q_{i,b,n}&=\frac{E_B(1-S_{i,b,n}^{+})}{\eta_i\delta},
\qquad Q_i=\frac{E_B}{\eta_i\delta}.
\end{aligned}
\label{lin:baseline_bounds}
\end{equation}
由 $0\le S_{i,b,n}^{+}\le1$ 可得 $0\le q_{i,b,n}\le Q_i$。若 $\bar p_{i,b}=0$，直接置 $P_{i,b,n}^{\mathrm{base}}=0$；若 $\bar p_{i,b}\ge Q_i$，直接置 $P_{i,b,n}^{\mathrm{base}}=q_{i,b,n}$。在 $0<\bar p_{i,b}<Q_i$ 时，引入 $\zeta_{i,b,n}\in\{0,1\}$，写成
\begin{equation}
\begin{aligned}
P_{i,b,n}^{\mathrm{base}}&\le\bar p_{i,b},\\
P_{i,b,n}^{\mathrm{base}}&\le q_{i,b,n},\\
P_{i,b,n}^{\mathrm{base}}&\ge\bar p_{i,b}\zeta_{i,b,n},\\
P_{i,b,n}^{\mathrm{base}}&\ge q_{i,b,n}-(Q_i-\bar p_{i,b})\zeta_{i,b,n}.
\end{aligned}
\label{lin:baseline_min}
\end{equation}
当 $\zeta_{i,b,n}=1$ 时功率取设备与站级额度共同确定的上限；当 $\zeta_{i,b,n}=0$ 时功率取本段可充电量对应的值。两种情况均满足两项上界，因而精确表示二者的最小值。运营结束条件与基准功率的乘积 $P_{i,b,n}=(1-d_n^{\mathrm{end}})P_{i,b,n}^{\mathrm{base}}$ 按式~\eqref{lin:product} 表达，其中基准功率的有效区间为 $[0,\min\{\bar p_{i,b},Q_i\}]$。

\begin{thebibliography}{9}
\bibitem{moreno2022terminal}
Moreno-Mora F, Beckenbach L, Streif S.
Predictive Control with Learning-Based Terminal Costs Using Approximate Value Iteration[EB/OL].
arXiv:2212.00361, 2022.
\url{https://arxiv.org/abs/2212.00361}.
\end{thebibliography}
\end{document}
